\documentclass[11pt,a4paper]{article}
\usepackage[T1]{fontenc}
\usepackage{lmodern}
\usepackage{amsmath,amssymb,amsthm,mathtools}
\usepackage[margin=27mm]{geometry}
\usepackage{microtype,booktabs,longtable,array,tabularx,enumitem}
\usepackage[hidelinks]{hyperref}
\setlength{\emergencystretch}{2em}
\setlength{\parskip}{3pt plus 1pt}
\setcounter{tocdepth}{1}
\newtheorem{theorem}{Theorem}[section]
\newtheorem{lemma}[theorem]{Lemma}
\newtheorem{proposition}[theorem]{Proposition}
\newtheorem{corollary}[theorem]{Corollary}
\theoremstyle{definition}
\newtheorem{definition}[theorem]{Definition}
\newtheorem{example}[theorem]{Example}
\theoremstyle{remark}
\newtheorem{remark}[theorem]{Remark}
\newcommand{\N}{\mathbb N}
\newcommand{\NN}{\mathbb N_{>0}}
\newcommand{\norm}{\operatorname{N}}
\newcommand{\lv}{\operatorname{lv}}
\newcommand{\NF}{\operatorname{NF}}
\newcommand{\sgn}{\operatorname{sign}}
\newcommand{\can}[1]{\mathopen{[\![}#1\mathclose{]\!]}}
\newcommand{\KK}[2]{K_{#1,#2}}
\newcommand{\EE}[2]{E_{#1,#2}}
\numberwithin{equation}{section}
\title{Identities and Comparison for the Standard Ordinal Collapse}
\author{}
\date{17 November 2026}
\hypersetup{pdftitle={Identities and Comparison for the Standard Ordinal Collapse}}
\begin{document}
\maketitle

\begin{abstract}
Fix the Wilken--Weiermann collapse below $\varepsilon_0$ with the
unchanged standard control $f(n)=F_5(n+100)$, where $F_0(n)=2^n$ and
$F_{j+1}(n)=F_j^{n+1}(n)$.  We study positive expressions obtained from
natural constants, positive natural inputs, addition, and the functions
$h_\lambda(u)=\psi(\lambda+u)$ for fixed nonzero limit heads.
Universal identities have complete recursive normal forms in any finite
number of inputs.  With one input, every nonidentity has an effectively
bounded finite equality set and a computable eventual sign.  For arbitrary
arrangements of polynomial heads, exact descent kernels and their complete
counts give a closed existential decision procedure.  Separate
parameter-preserving translations handle common-head and balanced layers
with arbitrary fixed heads; they use a published first-order arithmetic
theorem whose hypotheses are verified for this exact control.  A row
depending on one arbitrary positive core instead has a direct Presburger
description.  The corresponding source consequences include equality for
all ground-open canonical assignments, comparison with one ground free
name, all recursor-free sources, and all sources whose internal types have
level at most two.  These results concern the fixed standard assignment;
they complement the all-ground comparison theorem obtained by choosing a
different admissible control.
\end{abstract}
\tableofcontents

\section{Introduction}

An ordinal assignment can be effective to evaluate at a closed input and
still be difficult to compare at all open inputs.  In the assignment of
Wilken and Weiermann for G\"odel's $T$ \cite{ww}, the scalar coordinate
contains a norm-controlled collapse.  Its values grow by enormous jumps,
but this growth alone does not decide a signed comparison.  Two terms may
share a leading value and cancel it exactly, leaving a descent count or a
smaller correlated expression to determine the answer.

We fix the standard control throughout.  The distinction matters because
the definition permits other admissible controls.  The theorem that all
ground-open canonical comparisons are decidable for explicitly chosen
controls \cite{wwcomparison} does not replace the standard assignment by
an equal function.  The present results retain the original
$F_5(n+100)$ and address its finer algebraic and decision properties.

Three mechanisms will be used.  First, a collapse value is an exact sum
of a sparse terminal value and its complete descent count.  On suitable
inputs, distinct positive atoms have distinct terminal values; the largest
one separates their signed sum.  This gives free identity normal forms
and finite equality sets.  Second, retained-prefix kernels compare
arbitrary fixed ordinal heads in one input.  Their proof requires a
reserve large enough for an entire lower expression, and a comparison
algorithm whose separation certificates terminate with it.  Third,
finite-degree kernels admit exact counter expansions and constant
residues.  A ranked cancellation procedure then reduces a closed
multivariate problem to Presburger elimination, keeping all shared input
names and every finite exception.

Some multivariate expressions admit a simpler comparison.  At a common
head, input order also orders the sparse terminal values.  With different
heads, the same conclusion holds in a strict control interval.
Common-head and balanced certificates make these numerical facts into
formulas with the original free variables.  Unlike the general
finite-degree procedure, these translations can be placed under arbitrary
Boolean connectives and quantifiers.  Their final arithmetic decision is
supplied by Snyatkov's theorem \cite{snyatkov}; the exact hypotheses are
proved in Section~\ref{m:sec:control}.

\begin{center}
\small
\begin{tabularx}{\textwidth}{@{}>{\raggedright\arraybackslash}p{.27\textwidth}
 >{\raggedright\arraybackslash}X>{\raggedright\arraybackslash}X@{}}
\toprule
Input class & Effective conclusion & Essential restriction\\
\midrule
Positive expressions, any finite arity
& Complete universal identity normal form
& Fixed heads and positive inner syntax\\
Positive expressions, one input
& Finite equality set and eventual sign
& One shared raw input across the whole query\\
Arbitrary polynomial heads
& Closed existential and universal comparison
& One quantifier block; no claimed free-parameter elimination\\
Common-head or balanced layers
& Exact formulas and arbitrary quantified decisions
& One supplied finite certificate per comparison atom\\
One arbitrary positive core
& Finite Presburger sign formula
& All open dependence occurs through that core\\
\bottomrule
\end{tabularx}
\end{center}

The main results are statements about numerical expressions, so they do
not depend on every auxiliary expression having a source term.  The final
section applies them in the forward direction to literal canonical
assignments.  It states the source interface and distinguishes internal
type bounds from bounds on free variables.  All decision procedures are
mathematical termination results; no practical running-time bound is
asserted.

\section{Ordinal norms and effective collapse evaluation}
\label{f:sec:foundations}

The numerical assignment uses ordinals below $\varepsilon_0$, represented
by hereditary Cantor normal forms.  The ordinal $\varepsilon_0$ is the
least fixed point of $\alpha\mapsto\omega^\alpha$.  We distinguish
ordinary ordinal addition and multiplication from their natural
counterparts $\oplus$ and $\otimes$.  Natural addition collects the
Cantor monomials of its operands; natural multiplication distributes over
those monomials and naturally adds their exponents.  In particular both
natural operations are commutative.  All finite tails displayed as
$\lambda+n$ use ordinary addition.  If $\lambda$ is a limit and $n$ is
finite, then $\lambda\oplus n=\lambda+n$.

\begin{definition}[Additive hereditary norm]\label{f:def:additive-norm}
For decreasing exponents and positive integer coefficients put
\[
 N(0)=0,\qquad
 N\left(\sum_{i<r}\omega^{\alpha_i}c_i\right)
       =\sum_{i<r}c_i\bigl(1+N(\alpha_i)\bigr).
\]
Thus $N(n)=n$, $N(\omega)=2$ and $N(\omega^2)=3$.
Coefficients are charged with their full multiplicity.
\end{definition}

\begin{lemma}[Effective ordinal operations and norm bounds]
\label{f:lem:norm-arithmetic}
Every norm ball is finite and effectively enumerable.  Cantor-normal-form
comparison, $\oplus$, $\otimes$, and ordinal exponentiation $2^\alpha$
are effective.  Moreover
\begin{align*}
 N(\alpha\oplus\beta)&=N(\alpha)+N(\beta),\\
 N(\alpha)+N(\beta)-1
 &\le N(\alpha\otimes\beta)\le N(\alpha)N(\beta)
       &&(\alpha,\beta>0),\\
 N(\alpha)&\le2N(2^\alpha),\qquad
 N(2^\alpha)\le2^{N(\alpha)}.
\end{align*}
\end{lemma}
\begin{proof}
An exponent in a nonzero ordinal of norm at most $b$ has norm at most
$b-1$.  Induction on $b$ therefore enumerates the finite set of possible
exponents and the finite lists of coefficients whose total charge is at
most $b$.  Recursive lexicographic comparison decides their order.
The definitions of natural sum and product give finite algorithms.

For the norm bounds, expand coefficients into repeated monomials.
Suppose the two operands have respectively $r,s$ monomials, and let
$A,B$ be the sums of the norms of their exponents.  Their norms are
$r+A$ and $s+B$, while the product has norm $rs+sA+rB$.
Collecting equal monomials preserves these totals.  The upper product
bound follows by expanding $(r+A)(s+B)$; for the lower bound the
difference is
\[
 (r-1)(s-1)+(s-1)A+(r-1)B\ge0.
\]
Norm additivity for natural sum follows in the same way.

Write $\alpha=\omega\xi+n$, $n<\omega$, using ordinary
multiplication.  Then
\[
 2^\alpha=\omega^\xi 2^n,\qquad
 N(2^\alpha)=2^n(1+N(\xi)).
\]
Multiplication on the left by $\omega$ changes the exponent of each
monomial from $\gamma$ to $1+\gamma$.  Its norm stays the same or
increases by one, so
$N(\xi)\le N(\omega\xi)\le2N(\xi)$.
The last two asserted bounds now follow from
$1+a\le2^a$ for $a\in\N$ and
$2N(\xi)+n\le2^{n+1}(1+N(\xi))$.
These are the norm estimates of \cite[Proposition~2.2]{ww}, with the
additive convention made explicit.
\end{proof}

For reference, the finite hierarchy used throughout is
\[
 F_0(n)=2^n,\qquad F_{j+1}(n)=F_j^{\,n+1}(n),\qquad
 G(n)=F_5(n+100).
\]
We fix $f=G$ in this paper.  No residue correction or other change to the
control is made.  The estimates used below are proved in
Section~\ref{m:sec:control}.  Superscripts on a function denote finite
iteration, with the zeroth iterate equal to the identity.

\begin{definition}[Norm-controlled collapse]\label{f:def:collapse}
For $\alpha<\varepsilon_0$ set
\[
 \psi(\alpha)=\max\bigl(\{0\}\cup
   \{\psi(\beta)+1:\beta<\alpha,
                    \ N(\beta)\le f(N(\alpha))\}\bigr).
\]
Ordinal recursion defines this function, since each displayed maximum
is finite.
\end{definition}

\begin{lemma}[Basic collapse estimates]\label{f:lem:collapse-estimates}
If $\beta<\alpha$ and $N(\beta)\le f(N(\alpha))$, then
$\psi(\beta)+1\le\psi(\alpha)$.  For $n,k\in\N$,
\[
 \psi(n)=n,\qquad N(\alpha)\le\psi(\alpha),\qquad
 \psi(\alpha)+k\le\psi(\alpha+k).
\]
If $\alpha\ge\omega$, then $\psi(\alpha)>f(N(\alpha))$.
\end{lemma}
\begin{proof}
The first assertion is a term of the defining maximum.  Induction on
$n$ gives the finite values: every predecessor of $n$ is finite and
its norm is below $n\le f(n)$.  If $\alpha\ge\omega$, the finite
ordinal $f(N(\alpha))$ is an admitted predecessor, which proves the
strict estimate and hence the norm bound.  The finite case of the norm
bound has already been proved.  Finally $N(\alpha+k)=N(\alpha)+k$,
so $\alpha+k$ is an admitted predecessor of $\alpha+k+1$.
Induction on $k$ proves the finite-tail inequality.
\end{proof}

In particular a fixed-head collapse is increasing in its finite tail.
The lemma does not assert monotonicity of $\psi$ on arbitrary ordinal
arguments.  The norm premise must be verified whenever an ordinal
inequality is passed through the collapse.

\subsection{An exact maximum-predecessor algorithm}

For $\alpha>0$ and $b\in\N$ define
\[
 M(\alpha,b)=\max\{\beta<\alpha:N(\beta)\le b\},\qquad
 L(e,b)=M(\omega^e,b).
\]
The defining set for $M$ contains zero and is finite.  We also put
$M(0,b)=0$ as a convention; it is not a predecessor assertion at zero.

\begin{theorem}[Greedy maximum]\label{f:thm:maximum-predecessor}
The following simultaneous recursion computes $M$ and $L$ exactly.
Write $\alpha=\omega^e c+r$, where $c>0$ and $r<\omega^e$, and
put $w=1+N(e)$.
\begin{enumerate}[label=(\roman*)]
\item $M(\alpha,0)=0$ and $L(0,b)=L(e,0)=0$.
\item If $r>0$ and $cw\le b$, then
\[
 M(\alpha,b)=\omega^e c+M(r,b-cw).
\]
\item Otherwise set
\[
 d=\begin{cases}
 \min\{c-1,\lfloor b/w\rfloor\},&r=0,\\
 \min\{c,\lfloor b/w\rfloor\},&r>0.
 \end{cases}
\]
Then $M(\alpha,b)=\omega^e d+L(e,b-dw)$, with a zero block
omitted.
\item If $e,b>0$, put $a=M(e,b-1)$, $v=1+N(a)$ and
$(q,s)=\operatorname{divmod}(b,v)$.  Then
\[
 L(e,b)=\omega^a q+L(a,s).
\]
\end{enumerate}
\end{theorem}
\begin{proof}
If an affordable nonzero tail is present, retaining the entire leading
block and decreasing the tail dominates every predecessor with a smaller
leading block.  This proves (ii).  Otherwise the coefficient in (iii)
is the largest affordable coefficient compatible with the strict upper
bound.  Its remaining tail is less than $\omega^e$ and has precisely
the remaining norm budget.

For (iv), a nonzero candidate below $\omega^e$ with budget $b$ has
a leading exponent $a'<e$ satisfying $N(a')\le b-1$.
The largest possible leading exponent is therefore $a=M(e,b-1)$.
Cantor comparison then maximizes its coefficient and the smaller tail
in that order, yielding $q$ and $L(a,s)$.

Termination follows from the lexicographic measure (budget, phase),
with $M$ above $L$ at equal budget.  Clause (ii) lowers the budget;
(iii) either lowers it or passes from $M$ to $L$.  In (iv) the first
recursive call uses $b-1$, while $v\le b$ and $q\ge1$ give $s<b$
for the second.  This proves both effectiveness and exactness.
\end{proof}

\begin{lemma}[Limit saturation and monotonicity]\label{f:lem:limit-saturation}
For every nonzero limit $\lambda$ and every $b\in\N$,
\[
 N(M(\lambda,b))=b.
\]
For fixed $\lambda$ this maximum is strictly increasing in $b$.
For arbitrary bounds the function $M$ is weakly increasing in each
argument wherever the predecessor interpretation is defined.
\end{lemma}
\begin{proof}
If $\eta=M(\lambda,b)$ has unused budget $d=b-N(\eta)>0$,
then $\eta+d<\lambda$ and $N(\eta+d)=b$, contradicting
maximality.  For $b<c$, the predecessor $M(\lambda,b)+1$ is
below $\lambda$ and has norm at most $c$, proving strict increase.
The last assertion follows by inclusion of the defining candidate sets.
\end{proof}

\begin{theorem}[Single-predecessor collapse recurrence]
\label{f:thm:single-predecessor}
For every nonzero limit $\lambda$ and $n\in\N$,
\[
 \psi(\lambda+n)
   =\psi\bigl(M(\lambda,f(N(\lambda)+n))\bigr)+1.
\]
\end{theorem}
\begin{proof}
Induct on $n$.  Put $b=f(N(\lambda)+n)$ and
$\eta=M(\lambda,b)$.  It is an admitted predecessor, of norm $b$,
so it supplies the lower inequality.  Any other admitted $\beta<\lambda$
satisfies $\beta\le\eta$.  If $\beta<\eta$, then
$N(\beta)\le b=N(\eta)\le f(N(\eta))$; controlled strictness
therefore gives $\psi(\beta)+1\le\psi(\eta)$.

The remaining predecessors are $\lambda+j$, $j<n$.  For such a $j$
put $b_j=f(N(\lambda)+j)$ and $\eta_j=M(\lambda,b_j)$.
Strict growth of $f$ and limit saturation give
$\eta_j<\eta$ and $N(\eta_j)=b_j<b$.
The induction hypothesis and controlled strictness imply
\[
 \psi(\lambda+j)=\psi(\eta_j)+1\le\psi(\eta).
\]
Thus every contribution to the original maximum is at most
$\psi(\eta)+1$, which proves the reverse inequality.
\end{proof}

\begin{corollary}[Closed evaluation and its norm orbit]
\label{f:cor:closed-evaluation}
A closed collapse value is computable by the following deterministic
descent.  At a nonfinite state $\alpha=\lambda+n$, replace it by
$M(\lambda,f(N(\alpha)))$ and add one to an event counter.
On reaching a finite state $Z$, return $Z$ plus that counter.
If the initial norm is $t$ and the number of events is $D$, then
\[
 Z=f^D(t),\qquad \psi(\alpha)=Z+D.
\]
For a nonfinite initial state one also has $0<D<Z$.
\end{corollary}
\begin{proof}
Each step strictly lowers the ordinal, so well-foundedness below
$\varepsilon_0$ proves termination.  The preceding theorem preserves
the value after adding the event count, and finite evaluation gives the
last value.  Limit saturation makes the norm of each new state exactly
$f$ of the preceding norm, including the transition to the finite state.
There is no further event count along the finite predecessor chain.
A nonfinite start has $t\ge2$, and $f(n)>n$ gives
$Z\ge t+D>D$.
\end{proof}

\begin{corollary}[Finite coefficient identity]\label{f:cor:finite-coefficient}
For all $k,n\in\N$,
\[
 \psi(\omega k+n)=f^k(2k+n)+k.
\]
\end{corollary}
\begin{proof}
The case $k=0$ is finite evaluation.  If $k>0$ and
$b=f(2k+n)$, the exact maximum is
$M(\omega k,b)=\omega(k-1)+(b-2(k-1))$.
It has total norm $b$.  The single-predecessor recurrence and induction
on $k$ give the formula.
\end{proof}

The evaluator always calls $M$ with the \emph{limit part} of its
current ordinal.  Calling it with $\lambda+n$ instead would retain
finite predecessors that have already been accounted for in the maximum
argument and would not give the stated norm orbit.

\section{The standard control and its first-order arithmetic}
\label{m:sec:control}
Define, literally,
\begin{equation}\label{m:eq:f}
 F_0(n)=2^n,\qquad F_{j+1}(n)=F_j^{\,n+1}(n),\qquad
 f(n)=F_5(n+100).
\end{equation}
A superscript on a function denotes a finite iterate. The starting argument
in $F_{j+1}$ varies with $n$. This convention will not be replaced by
iteration from a fixed seed.

\begin{lemma}[Growth]\label{m:lem:growth}
Each $F_j$ is recursive and strictly increasing, with $F_j(n)\ge2^n$.
For $j\ge0$,
\[
 F_{j+1}(n+1)\ge2^{F_{j+1}(n)+1}.
\]
Consequently $f(z+1)\ge(A+2)f(z)$ whenever $z\ge A\ge0$.
\end{lemma}
\begin{proof}
Finite iteration preserves recursiveness, strict increase, and domination
of $2^n$. Increasing the input from $n$ to $n+1$ makes the first $n+1$
iterates at least one larger, and the definition then supplies one further
application of $F_j$. Thus
$F_{j+1}(n+1)\ge F_j(F_{j+1}(n)+1)\ge2^{F_{j+1}(n)+1}$.
For $v\ge1$, $2^{v+1}\ge(v+2)v$, by direct checking at $v=1,2,3$ and
induction. Apply this with $v=f(z)>z\ge A$.
\end{proof}

\begin{lemma}[Effective constant residues]\label{m:lem:residues}
For every $m\ge1$ one can compute $r_m\in[0,m)$ such that
\[
 F_j(n)\equiv r_m\pmod m\quad(j\ge1,\ n\ge m).
\]
In particular $f(n)\equiv r_m\pmod m$ for $n\ge\max(0,m-100)$.
For $s\ge1$, $j\ge2$, this can be strengthened to
\begin{equation}\label{m:eq:envelope}
 m\le 2^{2^s}\quad\Longrightarrow\quad F_j(s)\equiv r_m\pmod m.
\end{equation}
If $B=f(u)$ and $m\le2^{2^B}$, then also $f(u+1)\equiv r_m\pmod m$.
\end{lemma}
\begin{proof}
Set $r_1=0$. For $m>1$, write $m=2^a b$ with $b$ odd, and put
$q=\varphi(b)$ if $b>1$, $q=1$ otherwise. Since $q<m$, recursively
compute $r_q$, then use the Chinese remainder theorem to define
\[
 r_m\equiv0\pmod{2^a},\qquad r_m\equiv2^{r_q}\pmod b.
\]
The recursion strictly lowers the integer modulus.
For towers $T_0(n)=n$, $T_{h+1}(n)=2^{T_h(n)}$, induction on $m$ shows
$T_h(n)\equiv r_m\pmod m$ whenever $h,n\ge m$. Indeed, the last exponent
is at least $n\ge a$, while its residue modulo $q$ is $r_q$ by induction;
Euler's theorem handles the odd factor. This proves the assertion for
$F_1(n)=T_{n+1}(n)$. At every higher level, all intermediate arguments are
at least $n$, so the last application of the preceding level proves the
same residue assertion.

For the stronger bound, $F_1(s)\ge2^{2^s}$. In $F_2(s)$ there are at
least two $F_1$ applications; the final one has input at least $F_1(s)$,
so its input is at least $m$. The same argument propagates through every
$j\ge2$. Finally put $n=u+100$. The last $F_4$ call in
$F_5(n+1)=F_4^{n+2}(n+1)$ has input
$F_4^{n+1}(n+1)>F_4^{n+1}(n)=B$. Apply the stronger bound to that call.
\end{proof}

\begin{theorem}[Snyatkov's arithmetic supplier]\label{m:thm:snyatkov}
The first-order theory $\operatorname{Th}(\N;0,1,+,<,f)$, with the exact
function in \eqref{m:eq:f}, is decidable.
\end{theorem}
\begin{proof}
We use Theorem~2 and effective Corollary~2 of Snyatkov~\cite{snyatkov},
pp.~50--51. Its definitions on p.~40 require an increasing recursive
function effectively compatible with addition: effective eventual periodicity
modulo every integer, and an effective shifted sign separation for each
unbounded integer combination of fixed translates. It also requires effective
eventual periodicity of the \emph{fixed-seed} hyperfunction
$H(0)=1$, $H(n+1)=f(H(n))$. The cited theorem and effective corollary
give effective quantifier elimination after adjoining the specified
divisibility and floor-inverse symbols for $f$ and $H$.  Both functions
are recursive, strictly increasing and unbounded.  Their floor inverses
are computable by finite search, with value zero below the initial
function value.  The remaining atomic predicates are decidable.
Quantifier elimination therefore decides closed sentences in the
expanded language and in its $f$-only reduct.

The residue condition follows from Lemma~\ref{m:lem:residues}. For the second
condition combine equal shifts in $S(x)=\sum_b a_b f(x+b)$. If every
combined coefficient is zero the sum is zero. Otherwise let $B$ be the
largest shift with nonzero coefficient, and set
\[
 A=\sum_{b<B}|a_b|,\qquad
 \delta=\max(0,-\min b,A+1-B).
\]
Every argument at $x+\delta$ is nonnegative. With
$z=x+\delta+B-1\ge x+A$, Lemma~\ref{m:lem:growth} gives, if $a_B>0$,
\[
 S(x+\delta)\ge f(z+1)-Af(z)\ge2f(z)>f(x).
\]
For $a_B<0$, apply the same calculation to $-S$. This proves unboundedness,
its sign, and a suitable shift by finite arithmetic on the input data.
Lastly $H(n)\ge n+1$; for $n\ge m$, the last application of $f$ in
$H(n)$ has input at least $m$, so $H(n)\equiv r_m\pmod m$.
All published hypotheses are therefore paid. No identification of $H$ with
a diagonal $F_j$, and no decision result for variable-seed iteration, is
being imported.
\end{proof}


\section{Positive expressions, identities, and finite zeros}
\label{u:sec:positive}

\begin{definition}[Positive fixed-limit expressions]\label{u:def:positive}
For a finite list of variables $x_1,\ldots,x_d$ valued in $\NN$, the input
language $\mathcal P_d$ is generated by natural constants, those variables,
addition, and $h_\lambda(A)=\psi(\lambda+A)$, for fixed closed nonzero limits
$\lambda<\varepsilon_0$. Only finitely many heads occur in any input.
Natural coefficient multiplication means repeated addition. There is no
subtraction inside a head, multiplication of two open values, variable
ordinal index, inverse control or additional orbit-height operation.
\end{definition}
A difference $A-B$ is allowed as a signed row to be compared; it is not a
new constructor available inside an input expression. Indices may be
arbitrarily large below $\varepsilon_0$, and nesting is unrestricted.

Three conclusions will be established. Universal equality is decidable in
$\mathcal P_d$ for any finite $d$. In $\mathcal P_1$, a nonidentity has only
effectively bounded finitely many zeros. Finally, in $\mathcal P_1$ its
eventual sign is effective. The third conclusion requires substantially
more than the second: a nonzero integer-valued sequence can change sign
infinitely often without ever being zero.


The greedy evaluator of Corollary~\ref{f:cor:closed-evaluation} gives,
for an infinite ordinal input of norm $t$, the exact decomposition
\begin{equation}
  \psi(\alpha)=f^L(t)+L,\qquad L\ge1,                 \label{u:display:1-3}
\end{equation}
where $L$ counts replacements until the first finite state.  In particular
all displayed closed constants are computable.  The control estimates and
controlled strictness used throughout are
\begin{align}
 f(v)&\ge2^v>v,& f(v+1)&\ge2^{f(v)+1}, \label{u:eq:gaps}\\
 \alpha<\beta\ \&\ \norm(\alpha)\le f(\norm(\beta))
 &\ \Longrightarrow\ \psi(\alpha)<\psi(\beta).
 \label{u:eq:strict}
\end{align}
The first line follows from Lemma~\ref{m:lem:growth}; the second is a
direct clause of the collapse definition.  Neither asserts unrestricted
ordinal monotonicity of the collapse.

\begin{lemma}[Sparse peaks]\label{u:lem:sparse}
For an infinite input of norm $t\ge2$, write $Z=f^L(t)=f(s)$ as in \eqref{u:display:1-3}.
Then $0<L\le s-1$, and
\begin{equation}
           Z<\psi(\alpha)=Z+L<2Z<f(s+1).             \label{u:display:1-6}
\end{equation}
Distinct peaks $f(s)$ and $f(u)$ with $u<s$ satisfy $f(u)\le f(s-1)$.
\end{lemma}
\begin{proof}
Each norm update increases the norm by at least one. Thus
$s=f^{L-1}(t)\ge t+L-1\ge L+1$. Since $f(s)>s$, the count is smaller than
$Z$. The final strict inequality follows from \eqref{u:eq:gaps}. The last
assertion uses strict increase of the integer map $f$.
\end{proof}

\subsection{A complete normal form for identities}
\begin{definition}[Collected recursive normal form]\label{u:def:nf}
A normal form consists of a constant $c\in\N$, a vector
$b\in\N^d$, and a finite dictionary of tags $(\lambda,A)$ with positive
natural coefficients. Each $A$ is itself a normal form with nonempty
variable support. Its value is
\begin{equation}
       c+\sum_i b_i x_i+\sum_{(\lambda,A)} a_{\lambda,A}h_\lambda(A).
                                                               \label{u:display:2-1}
\end{equation}
The dictionary graph is finite and acyclic. The height is zero if no tag
occurs, and otherwise one plus the largest argument height.
\end{definition}
Normalize bottom up. Addition merges dictionaries and coefficient vectors.
For a head, normalize its argument; if that argument is closed, evaluate
the entire call and return a constant. Otherwise keep the tag
$(\lambda,\NF(A))$. Indices are compared as actual canonical Cantor forms.
A fixed finite tail in a presented index is first moved into its argument;
a zero limit part is removed using $\psi(n)=n$.

Identical normal forms have equal values, by induction. Completeness is a
separation theorem and needs more than commutative algebra.

\begin{lemma}[Orbit cancellation at roots]\label{u:lem:roots}
Call a natural number an $f$-root if it is outside $\operatorname{range}(f)$.
If infinite inputs $\alpha,\beta$ have root norms $t,u$ and equal peaks,
then $\alpha=\beta$.
\end{lemma}
\begin{proof}
From $f^L(t)=f^J(u)$ cancel the smaller number of iterates, using injectivity.
If $L\ne J$, one of the two roots is a positive iterate of the other, a
contradiction. Hence $L=J$ and $t=u$. Formula \eqref{u:display:1-3} gives equal collapse
values. If the ordinals differed, order them; their common norm satisfies
$t\le f(t)$, so \eqref{u:eq:strict} would give distinct values.
\end{proof}
This proof is useful because it avoids any assertion that terminal-peak
maps are monotone in arbitrary ordinal inputs.

\begin{theorem}[Identity normal form]\label{u:thm:identity}
For $E,F\in\mathcal P_d$, the identity $E=F$ holds at every positive
valuation exactly when $\NF(E)=\NF(F)$. If the forms differ, one can
effectively specify a positive countervaluation of the form
\begin{equation}
           x_i=w_i(f(k)+1),\qquad w_i,k\in\NN.        \label{u:display:2-2}
\end{equation}
Consequently universal equality in $\mathcal P_d$ is decidable.
\end{theorem}
\begin{proof}
If both expressions are closed, evaluate them; when they differ take
$k=1$ and every $w_i=1$ for the required countervaluation, with the
empty tuple allowed when $d=0$.  Otherwise let $\mathcal S$ contain
the two root normal forms, zero, and all recursively
reachable tag arguments. Let $M$ be the largest raw-variable coefficient
in this finite set, $B=\max(2,M+1)$, and $w_i=B^{i-1}$.
Uniqueness of finite base-$B$ notation shows that different coefficient
vectors in $\mathcal S$ have different weighted slopes $b\cdot w$.
On the ray $x_i=w_i n$, any open positive expression is at least $n$:
this follows by induction from positivity and $h_\lambda(q)\ge q$.

We first make every reachable open atom's initial norm a root on the ray
$n=f(k)+1$. For a tag $(\lambda,A)$ write that norm as
\begin{equation}
 t= A_0+Dn+\sum_j a_j h_{\lambda_j}(A_j),\quad
 A_0=\norm(\lambda)+c_A\ge2,\quad D=b_A\cdot w,
 \quad C=\sum_j a_j.                                 \label{u:display:2-3}
\end{equation}
If $C=0$, openness forces $D\ge1$. For $k\ge\max(4,A_0+2D)$,
\[
          f(k)<t\le(A_0+2D)f(k)<f(k+1),
\]
by \eqref{u:eq:gaps}. If $C>0$, let $Y=f(s)$ be the largest child peak.
Every child is open, so $Y\ge f(n)\ge2^n$. For
$n\ge\max(4,A_0,D+1,2C+1)$ the affine part is at most $n^2\le Y$.
The largest child contributes more than $Y$, while every child is below
$2Y$ by Lemma~\ref{u:lem:sparse}. Therefore
\begin{equation}
                 Y<t<(2C+1)Y<f(s+1).                \label{u:display:2-4}
\end{equation}
This too makes $t$ a root. All these bounds are finite syntax data.

Induct on the height of forms in $\mathcal S$. Distinct affine forms are
unequal for all sufficiently large $n$, because their slopes differ or
their constants differ. Suppose distinct lower-height forms have already
been separated. Two distinct tags at the next height cannot have equal
peaks: Lemma~\ref{u:lem:roots} would identify the ordinals
$\lambda+A(x)$ and $\mu+B(x)$, hence their limit parts and finite tails.
The indices would agree and the different lower-height arguments would
have equal values, contradicting induction.

Now cancel identical tags in the difference of two distinct forms at this
height. If tags remain, write the row on the weighted ray as
\begin{equation}
                \sum_j c_j h_{\lambda_j}(A_j)+pn+q,  \label{u:display:2-5}
\end{equation}
where $c_j\ne0$ and the remaining tags are distinct. Let $Y=f(s)$ be the
largest peak. Its coefficient has absolute value at least one. All other
peaks are at most $U=f(s-1)$, and their full values are below $2U$.
Because the largest tag is open, $s\ge n$, while $U\ge s$. Set
$K=2\sum_j|c_j|+|p|+|q|$. The absolute value of all lower contributions
is at most $KU$. Once $k\ge\max(4,K)$, the leading peak satisfies
$Y\ge2^{U+1}>KU$. Thus the row is nonzero, whichever tag is largest.
No algorithm for determining that largest tag has been assumed.

For explicit effectivity take the maximum of the root bounds above and,
for every pair in $\mathcal S$, the displayed row bound $K$; for an affine
row $pn+q$ with $p\ne0$, also include $|q|+1$. This single computable
$k$ suffices at every induction height. Normalization and all finite maxima
terminate, proving the theorem.
\end{proof}

\begin{example}[The exact boundary of positivity]
The collected forms of $h_\omega(x+y)+x$ and $x+h_\omega(y+x)$ agree.
In contrast the two tags in $h_\omega(x)+h_\omega(y)$ remain distinct;
the theorem decides an identity, not whether they happen to coincide at
some valuation. Inner subtraction cannot be added silently. The literal
finite-$\omega$ evaluation gives the exact identity
\[
 h_{2\omega}(n)=h_\omega(h_\omega(n+2)-3)+1
              =f(f(n+4))+2.
\]
The right side's subtraction is forbidden by Definition~\ref{u:def:positive}.
The free normal form is not asserted for that larger language.
\end{example}

\subsection{An effective finite-zero bound in one variable}
From now on every open input depends on the same $x\in\NN$.
This hypothesis permits a stronger separation statement along the entire
tail, without selecting the sparse ray \eqref{u:display:2-2}.

\begin{lemma}[Common-band or root separation]\label{u:lem:oneinput}
For any finite collection of normal forms in $\mathcal P_1$, a computable
$N$ has the following properties for every $x\ge N$.
Every affine open atom input lies in $[x,f(x))$. Every nonlinear atom input
is a root strictly greater than $f(x)$. Equal peaks of two tags force equal
indices and equal argument values.
\end{lemma}
\begin{proof}
Write an atom norm as $A_0+bx+\sum a_jh_{\lambda_j}(A_j)$ as in \eqref{u:display:2-3}.
For an affine input, $b\ge1$ and
$x\le A_0+bx<f(x)$ once $x\ge\max(4,A_0+b+1)$.
For a nonlinear input the largest child peak $Y$ is at least $f(x)$;
for $x\ge\max(4,A_0+b,2C+1)$ the argument proving \eqref{u:display:2-4} applies
at every such $x$. It is a root greater than $f(x)$.

Two distinct inputs in $[x,f(x))$ cannot lie on the same $f$-orbit,
since the first positive iterate of the smaller is already at least
$f(x)$. Two nonlinear roots are handled by cancellation. In a mixed pair,
orbit cancellation either puts the larger nonlinear root in the range of
$f$, expresses the smaller affine input as an iterate of that larger
input, or identifies their disjoint bands; each is impossible. Thus equal
peaks force equal initial norms and equal iterate counts. Equation \eqref{u:display:1-3}
and controlled strictness then identify the two ordinals, exactly as in
Lemma~\ref{u:lem:roots}.
\end{proof}

\begin{theorem}[Effective finite zeros]\label{u:thm:zeros}
Given $E,F\in\mathcal P_1$, one computes $N\ge1$ such that either their
normal forms agree, or $E(x)\ne F(x)$ for every $x\ge N$. In the latter
case their equality set is the computable finite set obtained by exact
evaluation at $1,\ldots,N-1$.
\end{theorem}
\begin{proof}
Use the same finite set $\mathcal S$ as in Theorem~\ref{u:thm:identity}.
Take all bounds from Lemma~\ref{u:lem:oneinput}. For each pair in
$\mathcal S$, cancel its tag dictionaries. If a tag remains, include
$4$ and $2\sum|c_j|+|p|+|q|$ in the maximum; otherwise include
$|q|+1$ when its affine slope $p$ is nonzero. Height induction now
separates distinct argument forms and, by Lemma~\ref{u:lem:oneinput}, the
peaks of distinct tags. The sparse leading-peak estimate \eqref{u:display:2-5} separates
the full rows on this entire tail. These finitely many bounds are computed
before the induction and work simultaneously at every height.
\end{proof}

\begin{lemma}[Range cells and signed rows]\label{u:lem:cells}
The following assertions have effective tail bounds.
For a nonlinear positive expression $E$, let $Z_i$ be its displayed open
tag peaks, and $Y=\max_i Z_i=f(s)$. Then
\begin{equation}
              Y<E<f(s+1),\qquad E<h_\omega(Y).        \label{u:display:3-1}
\end{equation}
For a value $V\in\operatorname{range}(f)$ on that tail,
\begin{equation}
 V<E\ \Longleftrightarrow\ (\exists i\ Z_i\ge V),\qquad
 V>E\ \Longleftrightarrow\ (\forall i\ Z_i<V).       \label{u:display:3-2}
\end{equation}
After cancellation of identical tags in a signed row, its eventual
pointwise sign is the coefficient sign of its greatest remaining peak,
once those peaks have been ordered. If no tags remain, the row is affine.
\end{lemma}
\begin{proof}
Every open child uses the same input $x$, so its peak dominates the entire
fixed affine part on an effective tail. The proof of \eqref{u:display:2-4}, with the
outer index norm omitted, gives $Y<E<(2C+1)Y<f(s+1)$.
Increasing the threshold makes this fixed multiple smaller than
$h_\omega(Y)=f(Y+2)+1$. There is no range value strictly between $f(s)$
and $f(s+1)$, proving \eqref{u:display:3-2}, including impossibility of $V=E$.
Theorem~\ref{u:thm:zeros}'s height induction separates the remaining tag peaks,
and the estimate in \eqref{u:display:2-5} proves the signed-row assertion. It does not
identify the greatest peak; that task is supplied below.
\end{proof}


\section{Principal kernels, retained mass, and lower profiles}
\label{u:sec:kernels}
The sign algorithm uses auxiliary kernels to expose the terminal peak of
a fixed source head. Their indices and masses are closed finite data,
and auxiliary kernels are not new input constructors in $\mathcal P_1$.

\begin{definition}[Principal suffix kernel]\label{u:def:kernel}
Let $\pi=\omega^\beta\ge\omega$ be a fixed principal, $a\in\N$, and
$t\ge a+\norm(\pi)$. Start at suffix
$\pi+(t-a-\norm(\pi))$, with total norm $t$ and permanent numerical mass
$a$. At a nonfinite suffix $\eta+n$, replace it by
$M(\eta,f(v)-a)$, where $v$ is the current total norm.
Let $\EE\pi a(t)$ count steps to the first finite suffix and put
\begin{equation}
              \KK\pi a(t)=f^{\EE\pi a(t)}(t).        \label{u:display:4-1}
\end{equation}
The finite exit suffix is $\KK\pi a(t)-a$. The value in \eqref{u:display:4-1} is the
\emph{total} exit norm, and always belongs to $\operatorname{range}(f)$.
\end{definition}
The suffix is strictly descending and its next norm budget is always
genuine. Thus the procedure terminates. A numerical mass need not be the
norm of any source prefix.

\begin{lemma}[Exact prefix splitting]\label{u:lem:prefix}
Let $\kappa$ be a fixed Cantor prefix and let $\eta$ be a nonzero limit
whose leading monomial is at most the least monomial of $\kappa$;
the zero prefix is also allowed. In particular one may take
$\kappa=\pi k$ and $\operatorname{lead}(\eta)\le\pi$.
For $b\ge c=\norm(\kappa)$,
\begin{equation}
 M(\kappa+\eta,b)=\kappa+M(\eta,b-c).                \label{u:display:4-2}
\end{equation}
A descent above $\kappa$ reaches a state $\kappa+n$ with finite remainder
before passing below $\kappa$. During that segment the remainder's exact
norm control is
\begin{equation}
                   G_c(v)=f(c+v)-c\ge f(v)>v.         \label{u:display:4-3}
\end{equation}
\end{lemma}
\begin{proof}
The ordinal $\kappa$ is eligible below $\kappa+\eta$, so the maximum cannot
have a smaller prefix. Every ordinal in the interval has a unique remainder
$r<\eta$, and compatible Cantor concatenation gives norm $c+\norm(r)$.
Maximizing that remainder proves \eqref{u:display:4-2}. Equal leading boundary exponents
cause coefficients and their norm contributions to add, so they are allowed.
Budgets increase, keeping the prefix affordable. Equation \eqref{u:display:4-2} forbids
skipping the finite-remainder cut; ordinal well-foundedness supplies it.
Norm saturation proves \eqref{u:display:4-3}. An increasing integer map preserves integer
differences, giving its inequality.
\end{proof}

\begin{lemma}[Ordered controls and counts]\label{u:lem:counts}
For increasing integer controls $g\le h$, both strictly exceeding their
input, let $L_g,L_h$ be the respective norm-maximal descent counts.
Then
\begin{equation}
 \alpha\le\beta,\quad\norm(\alpha)\le\norm(\beta)
       \quad\Longrightarrow\quad L_g(\alpha)\le L_h(\beta). \label{u:display:4-4}
\end{equation}
\end{lemma}
\begin{proof}
Induct on $\beta$. Finite $\alpha$ has zero count. Otherwise write the
limit parts of $\alpha,\beta$ as $\lambda,\mu$. If $\lambda=\mu$, their
next states $p,q$ are ordered by monotonicity in the budget, and
$\norm(p)=g(\norm\alpha)\le h(\norm\beta)=\norm(q)$.
If $\lambda<\mu$, then $\alpha<\mu$ and its norm is less than the next
$h$-budget. Hence $\alpha$ is eligible below $\mu$, so
$p<\alpha\le q$; the same norm inequality holds. Apply induction at
$q<\beta$ and add the first step.
\end{proof}
This also proves fixed-head monotonicity: increasing the finite tail of
$\rho+n$ increases its norm and does not decrease its count, so its full
$\psi$ value is strictly increasing. Fixed-mass kernels strictly increase
their genuine input by the same argument. Decreasing their retained mass
at fixed total input weakly increases the suffix and its budgets, so cannot
decrease the exit norm.

\begin{lemma}[Exact peak word]\label{u:lem:peakword}
For each fixed limit $\lambda$, one computes a finite word $W_\lambda$ of
principal kernels such that the peak of $h_\lambda(A)$ is exactly
\begin{equation}
             Z_{\lambda,A}=W_\lambda(\norm(\lambda)+A). \label{u:display:4-5}
\end{equation}
The omitted count is positive: $h_\lambda(A)=Z_{\lambda,A}+D_{\lambda,A}$.
\end{lemma}
\begin{proof}
Expand $\lambda$ into coefficient-one principal Cantor blocks. Delete them
from right to left. For each block $\pi$, the remaining earlier prefix has
a fixed norm $a$. Lemma~\ref{u:lem:prefix} says that this segment changes the
current total norm by exactly $\KK\pi a$. The last segment has prefix
zero. Concatenating the finitely many segments gives the exact word.
Every removed nonfinite segment contributes its literal count to $D$;
no count is set to zero in \eqref{u:display:4-5}.
\end{proof}

\subsection{Paying an entire lower positive profile}
A lower profile below a principal $L>\omega$ means a fixed expression
$P(z)$ formed from constants, $z$, addition and heads $h_\rho$ with
$\rho<L$. Its heads may themselves be nonprincipal limits.
The following reserve argument is needed to compare unequal kernel inputs
and to handle a limit exponent. Merely knowing that a larger kernel exceeds
a fixed number of $f$-iterates would be insufficient.

\begin{lemma}[One-stage reserve consumer]\label{u:lem:stage}
Let $\rho<\pi m$ be a fixed nonzero limit, $m\ge1$, and begin an ordinary
descent at $A=\pi(k+m)+n$, of norm $X\ge x\ge\norm(\rho)$.
At its first cut $\pi k+n'$ with finite remainder,
\begin{equation}
        n'>J_\rho(x):=f^{D_\rho(x)}(x)+D_\rho(x),    \label{u:display:5-1}
\end{equation}
where $D_\rho(x)$ is the count from
$\rho+(x-\norm(\rho))$.
\end{lemma}
\begin{proof}
Put $c=k\norm(\pi)$ and $\alpha=\rho+(x-\norm\rho)$.
The ordinal $\pi k+\alpha$ is below the initial limit part and has norm
$c+x\le2X<f(X)$, hence is eligible for the first predecessor. Its
remainder is at least $\alpha$ and has norm $f(X)-c\ge x$.
Lemmas~\ref{u:lem:prefix} and \ref{u:lem:counts} show that the following
prefix segment takes at least $d=D_\rho(x)$ further steps. Its exit total
norm $Y$ is at least $f^{d+1}(X)$. As $c\le X$ and every increasing
integer iterate preserves differences,
\[
 n'=Y-c\ge f^{d+1}(X)-X\ge f^{d+1}(x)-x.
\]
Write $W=f^d(x)\ge x+d$. Then
$f(W)-x>2W-x\ge W+d=J_\rho(x)$, as required.
Notice that the prefix mass was paid using $c\le X$; it need not be
bounded by the smaller test input $x$.
\end{proof}

\begin{lemma}[Effective profile reserve]\label{u:lem:reserve}
From a lower profile $P$ below $L>\omega$ one computes a nonempty word
\begin{equation}
 F=h_{\rho_l}\circ\cdots\circ h_{\rho_1},\quad
 F(z)>P(z)\ (z\ge z_0),\quad
 R=\sum_j(\norm(\rho_j)+1),\quad
 H=\max_j\norm(\operatorname{lead}(\rho_j)),          \label{u:display:5-2}
\end{equation}
with every $\rho_j<L$. If $\pi\ge\omega$ is principal and at least every
displayed leading monomial, a descent from $\pi R+n$, $n\ge z\ge z_0$,
has terminal norm greater than $P(z)$. The same conclusion holds with any
additional permanent numerical mass $c$.
\end{lemma}
\begin{proof}
Choose $z_0\ge4$, all constants and every weighted addition fan-in in $P$.
Process its expression graph in topological order, maintaining a common
upper bound $u$ on all node values already visited. An addition with fan-in
$b$ is bounded by $bu\le u^2\le2^u<h_\omega(u)$; append $h_\omega$ to the
running bound. For a head node append that head, using its monotonicity.
Every update increases the running bound. Append a final $h_\omega$ to
make the bound strict and the word nonempty. Since $L>\omega$, all added
heads are still below $L$.

For a head $\rho_j$ take $m_j=\norm(\rho_j)+1$ blocks. Its leading
coefficient is at most its norm, so $\rho_j<\pi m_j$. Moreover the
current norm at $\pi(k+m_j)+n$ is at least $\norm(\rho_j)+n$.
Lemma~\ref{u:lem:stage}, with $x=\norm(\rho_j)+n$, therefore makes the
finite tail at the next cut exceed $h_{\rho_j}(n)$. Repeat through the
word, using exactly $R$ blocks. The final finite norm exceeds $F(z)>P(z)$.
For permanent mass $c$, the exact suffix control is $G_c\ge f$.
Lemma~\ref{u:lem:counts} orders the counts and the iterated norms, so the
terminal suffix norm is at least the unprefixed one; adding $c$ only
increases the total norm.
\end{proof}

\begin{lemma}[Direct kernel growth]\label{u:lem:direct}
For fixed $L>\omega$ principal and fixed lower profile $P$, one computes
$T$ such that, for $t\ge T$ and every genuine mass $0\le c\le t-\norm(L)$,
\begin{equation}
                         \KK Lc(t)>P(t).             \label{u:display:5-3}
\end{equation}
Also, $\KK\pi c(t)\le h_\pi(t)$ on its genuine domain.
\end{lemma}
\begin{proof}
Choose one fixed principal $\omega\le\pi<L$ dominating the finite list of
leading monomials in \eqref{u:display:5-2}. For
$t\ge\max(z_0,R\norm\pi,4,\norm L)$, the ordinal $S=\pi R+t<L$ has
\[
             c+\norm(S)\le R\norm\pi+2t\le3t<f(t).
\]
It is eligible below the actual first suffix $M(L,f(t)-c)$. Count
comparison under the same induced control transfers the reserve bound
from $S$ to that maximum. This proves \eqref{u:display:5-3} uniformly in the masses.
For the final assertion first lower the retained mass to zero, then
increase the seed from $t$ to $\norm\pi+t$. Mass monotonicity and
fixed-mass input monotonicity bound the kernel peak by the full head value.
\end{proof}

\subsection{Adjacent inputs and limit-exponent mass separation}
\begin{theorem}[Cross-mass strict-input separation]\label{u:thm:adjacent}
For a fixed principal $L>\omega$ and fixed lower profile $P$ with reserve
data \eqref{u:display:5-2}, put
\begin{equation}
              T=\max(4,R+5,z_0,\norm L,H).            \label{u:display:6-1}
\end{equation}
For $t\ge T$, $0\le c\le t+1-\norm L$, and
$0\le e\le t-\norm L$,
\begin{equation}
       \KK Lc(t+1)>P(\KK Le(t)+\EE Le(t)).            \label{u:display:6-2}
\end{equation}
There is no ordering requirement on $c,e$. More generally the same
conclusion holds with any $u>t$ and any mass genuine at $u$ on the left.
\end{theorem}
\begin{proof}
Write $B=f(t)$, $B'=f(t+1)$, $b=M(L,B-e)$,
$\pi=\operatorname{lead}(b)$ and $w=\norm\pi$. Since
$B-e\ge f(t)-t+\norm L\ge t+\norm L\ge H$, every leading monomial
needed by \eqref{u:display:5-2} is eligible below $L$ in this budget. Also $B-e\ge2$,
so $\pi\ge\omega$. Thus $\pi<L$, $w\le B$, and $\pi$ dominates
every required leading monomial.

Put $p=c+Rw$ and
\[
 k=B'-B-(p-e),\qquad A=\pi R+b+k.
\]
Here $p-e\le t+1+RB\le(R+1)B$ and
$B'\ge2^{B+1}>(R+2)B$, so $k\ge0$. The sum is compatible CNF,
has norm $B'-c$, and lies below the principal $L$. Hence
$A\le Q=M(L,B'-c)$ and both have the same norm. The latter is the
upper process's actual first suffix.

Follow the packed suffix $A$ retaining its first $R$ blocks while deleting
the remainder $b+k$. At the reserve cut it has finite tail
\begin{equation}
               n=V-p,\qquad V=f^q(B').              \label{u:display:6-3}
\end{equation}
The lower process, after its first step, starts at $b$ with suffix norm $B-e$ and total norm $B$
and retained mass $e$. Let its remaining count be $\ell$, and set
$D=1+\ell=\EE Le(t)$ and $Z=f^\ell(B)=\KK Le(t)$.
The packed remainder initially exceeds $b$. The initial global gap pays
the mass difference $p-e$ and, as increasing integer iterates preserve
differences, at every simultaneous later step
$f^i(B')-p\ge f^i(B)-e$. Monotonicity of $M$ in its limit bound and
budget keeps the packed remainder above the lower remainder; it cannot
become finite earlier. Hence $q\ge\ell$.

If $s=f^{D-1}(t)$, then
\[
 V\ge f^\ell(B')=f^D(t+1)\ge f(s+1)\ge2^{Z+1}.
\]
Furthermore $D<Z$ and $p\le(R+1)Z$. Since $Z\ge t\ge R+5$,
$p+Z+D<(R+3)Z<Z^2<2^{Z+1}\le V$. Thus $n>Z+D$ after the full
retained mass has been subtracted. Lemma~\ref{u:lem:reserve} gives terminal
total norm greater than $P(Z+D)$ for the packed process. Count comparison
from $A\le Q$ at their equal norms, under their common control $G_c$,
transfers that bound to the actual upper kernel.

For arbitrary $u>t$, use the worst-mass inequalities
\begin{equation}
 \KK Lc(u)\ge \KK L{u-\norm L}(u)
          \ge \KK L{t+1-\norm L}(t+1).              \label{u:display:6-4}
\end{equation}
The first is mass monotonicity. For the second both processes start at
the identical suffix $L$, while their simultaneous suffix budgets are
$f^i(u)-u+\norm L\ge f^i(t+1)-(t+1)+\norm L$.
Their counts and terminal total norms are ordered. Apply the adjacent
case with the displayed genuine upper mass.
\end{proof}

At $L=\omega$ the kernel is simply $f(t)$, independently of mass. There
are no nonzero lower limit heads, so $P(z)=az+b$. Taking
$T_\omega=\max(4,a+b+2)$ gives
$P(f(t)+1)<f(t+1)$ by \eqref{u:eq:gaps}. This is the separate base case of
Theorem~\ref{u:thm:adjacent}.

\begin{theorem}[Limit-exponent mass separation]\label{u:thm:limitmass}
Let $L=\omega^\beta$ with $\beta$ a nonzero limit. For a fixed lower
profile $P$ and data \eqref{u:display:5-2}, put
\begin{equation}
           T=\max(4,R+4,z_0,\norm L+1,H).            \label{u:display:6-5}
\end{equation}
For $t\ge T$ and $0\le c<e\le t-\norm L$,
\begin{equation}
                \KK Lc(t)>P(\KK Le(t)+\EE Le(t)).    \label{u:display:6-6}
\end{equation}
Consequently the lexicographically larger pair $(\text{input},-\text{mass})$
gives domination against every fixed lower profile, on an effective
genuine tail. Equal pairs give equal values.
\end{theorem}
\begin{proof}
Put $B=f(t)$ and $C=f(B)$. Because $\beta$ is a limit, the maximum
affordable exponent $\delta_a=M(\beta,B-a-1)$ saturates its norm budget.
Its monomial $\pi_a=\omega^{\delta_a}$ already costs $B-a$, so the first
suffix below $L$ is exactly $\pi_a$, with coefficient one and no remainder.
Since $c<e$, saturation gives $\delta_c>\delta_e$ and $\pi_c>\pi_e$.
Put $w=\norm\pi_e=B-e$. The inequality
$B-e\ge f(t)-t+\norm L\ge t+\norm L\ge H$ shows that $\pi_e$ is
infinite and dominates every required leading monomial.

At the second upper step pack
\[
       k=C-c-(R+1)w,\qquad A=\pi_e(R+1)+k.
\]
Since $C>(R+1)B$ and $c+(R+1)w\le(R+1)B$, the tail $k$ is nonnegative.
The ordinal $A$ is below $\pi_c$ and has norm $C-c$, hence is below the
actual second upper suffix $Q=M(\pi_c,C-c)$ at that same norm.

In the packed process hold $R$ blocks while deleting $\pi_e+k$.
The retained mass is $p=c+Rw$. Compare this remainder with the lower
process after its first step, namely $\pi_e$ with total norm $B$ and
retained mass $e$. The global difference $C-B$ exceeds $p-e\le RB$,
so synchronized budgets keep the packed remainder above the lower one.
If the latter needs $\ell$ further steps, the reserve cut has norm
$V\ge f^\ell(C)=f(Z)$, where $Z=\KK Le(t)=f^\ell(B)$ and
$D=\EE Le(t)=1+\ell<Z$. Since $p\le(R+1)Z$ and $Z\ge R+4$,
$f(Z)>(R+3)Z>p+Z+D$. Its finite tail after all retained mass is paid
is therefore greater than $Z+D$.

Consume the reserve by Lemma~\ref{u:lem:reserve}, still with permanent mass
$c$. Transfer from $A\le Q$ at equal suffix norms using
Lemma~\ref{u:lem:counts}. This yields \eqref{u:display:6-6}. Every estimate uses only the
displayed genuine domains, not a fixed bound on $c,e$, so it is uniform
over the stated numerical masses. Combine with
Theorem~\ref{u:thm:adjacent} for unequal inputs to obtain the lexicographic
assertion. A profile-dependent threshold is intended; there is no single
threshold valid for all profiles at once.
\end{proof}
The limit-exponent hypothesis is essential. For example, the successor
identity proved next gives
$\KK{\omega^2}0(t)=f(\KK{\omega^2}2(t))$, while the lower profile
$h_\omega(z)=f(z+2)+1$ exceeds $f(z)$. Thus arbitrary lower-profile
absorption at the same input is false for successor exponents.

\subsection{Successor cuts and all same-input mass classes}
\begin{lemma}[Exact successor cut and count cocycle]\label{u:lem:successor}
Let $\theta=\omega^\gamma\ge\omega$, $L=\theta\omega$, and
$w=\norm\theta$. On the common genuine domain $t\ge c+w+\norm L$,
\begin{align}
 \KK Lc(t)&=\KK\theta c(\KK L{c+w}(t)),              \label{u:eq:cut}\\
 \EE Lc(t)&=\EE L{c+w}(t)
        +\EE\theta c(\KK L{c+w}(t)).                \label{u:eq:cocycle}
\end{align}
\end{lemma}
\begin{proof}
For $b=wq+r$, $0\le r<w$, the literal maximum is
\begin{equation}
                   M(L,b)=\theta q+M(\theta,r).      \label{u:display:7-3}
\end{equation}
The largest monomial below $L$ is $\theta$, costing $w$; greedily maximize
its coefficient, then the remainder below $\theta$. This also covers
$q=0$. Lowering the first budget from $f(t)-c$ to $f(t)-c-w$ removes
exactly one $\theta$ block and leaves the same remainder. Retaining that
one block in the $c$-process makes its mass during deletion of the other
remainder exactly $c+w$. Both processes have total norm $f(t)$ after
their first step. Prefix splitting therefore gives identical subsequent
norms until the latter process's finite cut. The remaining block is
exactly $\KK\theta c$, and the counts add. The stronger common domain
ensures that the inner shifted-mass kernel exists.
\end{proof}

\begin{lemma}[Small principal kernels]\label{u:lem:small}
On their genuine domains,
\begin{equation}
 \KK\omega c(t)=f(t),\quad\EE\omega c(t)=1,\qquad
 \EE{\omega^2}c(t)=f(t)/2-\lceil c/2\rceil+1.         \label{u:display:7-4}
\end{equation}
Thus $\omega^2$ kernels at the same input have exactly the mass classes
$\lceil c/2\rceil$, ordered in reverse. The control has even values.
\end{lemma}
\begin{proof}
Below $\omega$ the first predecessor is finite. Below $\omega^2$ it is
$\omega q+r$ with $q=\lfloor(f(t)-c)/2\rfloor$ and $0\le r<2$.
The first step is followed by exactly one step for each $\omega$ block,
because each such block's deletion uses a finite suffix maximum. Thus
the count is $q+1$, giving \eqref{u:display:7-4}. The definition of the iterated
exponential hierarchy makes each $f(t)$ even.
\end{proof}

\begin{theorem}[Strict mass order above degree two]\label{u:thm:mass}
For every fixed principal $L\ge\omega^3$, a computable $T_L$ satisfies
\begin{equation}
 t\ge T_L,\quad0\le c<e\le t-\norm L
                 \quad\Longrightarrow\quad\KK Lc(t)>\KK Le(t). \label{u:display:7-5}
\end{equation}
For a nonzero limit exponent take $T_L=\max(7,\norm L+1)$.
For $L=\theta\omega$, $\theta\ge\omega^2$, $w=\norm\theta$, take
\begin{equation}
 T_L=\max(4,\norm L,2w,\max(8,w),S_\theta),          \label{u:display:7-6}
\end{equation}
where $S_\theta=0$ at $\theta=\omega^2$ and otherwise $S_\theta=T_\theta$.
\end{theorem}
\begin{proof}
The limit case is Theorem~\ref{u:thm:limitmass} with profile $P(z)=z$.
Its compiler uses just one final $h_\omega$, giving $R=3,H=2,z_0=4$.
For the successor case let $B=f(t)$, and write the first suffixes from
\eqref{u:display:7-3} as $\theta q_a+r_a$, $a=c,e$. Here $q_c\ge q_e=q$ and
$r_a=M(\theta,(B-a)\bmod w)$. The domain and \eqref{u:display:7-6} imply
$B-e\ge f(t)-t+\norm L\ge2w$, so $q\ge2$.

Hold $q$ common $\theta$ blocks while deleting residuals
$\theta(q_c-q)+r_c$ and $r_e$. The first residual is at least the
second. Their permanent masses are $c+wq$ and $e+wq$; at equal global
norms their suffix budgets preserve this order. The finite-cut entry
norms therefore satisfy $U_c\ge U_e\ge B$.

With $k$ common blocks remaining, the child masses are
$a_k=c+w(k-1)<b_k=e+w(k-1)$. They are genuine even at the smaller
entry norm $V\ge B$, because
\begin{equation}
 b_k\le e+w(q-1)=B-((B-e)\bmod w)-w\le B-w\le V-w. \label{u:display:7-7}
\end{equation}
If the entry norms are unequal, Theorem~\ref{u:thm:adjacent} with identity
profile preserves strict order, uniformly over both masses; its threshold
is $\max(8,w)$. If they agree and $\theta\ge\omega^3$, the inductive
strict-mass theorem makes the first child outputs strictly ordered, and
every subsequent block preserves it.

For $\theta=\omega^2$, Lemma~\ref{u:lem:small} gives weak order, with a
tie precisely when $b_k=a_k+1$ and $a_k$ is odd. If the first block ties,
the next pair of masses subtracts $w=3$ from both; the smaller mass is
then even, so their ceilings differ and the second outputs are strictly
ordered. There are at least two blocks, and \eqref{u:display:7-7} pays their domains.
All further blocks preserve the strict order. This proves the base and
the successor induction. Removing the finite terminal successor tail of
the exponent reaches either this base or a limit exponent, so the
threshold recursion terminates.
\end{proof}


\section{The complete one-variable comparison algorithm}
We now supply the order decision that the normal-form proof deliberately
did not assume. There are two sorts in the auxiliary comparison language.
The source sort $P$ is $\mathcal P_1$. A range expression $R$ is a
nonempty finite word of kernels applied to one open $P$ expression.
Every argument of a kernel is either $P$ or $R$; a sum of range expressions
is not introduced into a kernel argument. Closed words are evaluated and
folded into constants. Every open expression of either sort is at least
$x$ on a computable tail. All generated masses and indices are fixed
natural or ordinal data on each branch.

\begin{theorem}[Effective eventual signs]\label{u:thm:sign}
Given two expressions in these sorts, one computes a sign
$\sigma\in\{-1,0,1\}$ and $N\ge1$ such that their difference has sign
$\sigma$ for every $x\ge N$. In particular this holds for every pair in
$\mathcal P_1$.
\end{theorem}
The proof is a simultaneous comparison and separation-certificate algorithm.
We specify its rules before proving their common termination measure.
Every rule retains the thresholds of all its children and the finitely
many kernel domains it uses. Since every open argument is at least $x$
eventually, each fixed genuine-domain obligation is met by increasing
$N$ to a computable constant.

\subsection{The comparison rules}
For $P/P$, form the signed normal-form row. An affine row is decided
directly. If tags remain, replace each relevant tag by its exact peak
word \eqref{u:display:4-5}, compare those $R$ expressions recursively and choose the
greatest peak. Distinct tags cannot have equal peaks beyond the bound of
Theorem~\ref{u:thm:zeros}. Return its coefficient sign, using
Lemma~\ref{u:lem:cells}. Identical forms return exact equality.

For $R/P$, a constant or nonconstant affine $P$ is eventually smaller:
the open $R$ is at least $f(x)$, which dominates every fixed affine expression.
For a nonlinear $P$ compare $R$ with its displayed peaks \eqref{u:display:4-5}, and use
the exact range-cell test \eqref{u:display:3-2}. The $P/R$ rule is symmetric.

For $R/R$ with equal outer principal, compare the two literal arguments
recursively. Strictly ordered arguments give the same strict kernel order
by Theorem~\ref{u:thm:adjacent}, or by monotonicity of $f$ at $\omega$.
If the arguments agree eventually, use the following same-input key:
mass is irrelevant at $\omega$; compare $\lceil c/2\rceil$ in reverse at
$\omega^2$; compare masses in reverse above $\omega^2$ by
Theorem~\ref{u:thm:mass}. Equal keys give equal outer values.

For different outer principals $H>M$, if $H$ has a successor exponent,
use the exact cut \eqref{u:eq:cut} on the $H$ side until the new outer
principal is $M$ or has a limit exponent still above $M$.
This cannot jump below $M$: every principal below $\theta\omega$ is at
most $\theta$. It takes only the finite terminal successor part of the
exponent. Combine these rewrites immediately with the equal-principal
rule or the following mixed-limit rule; do not recursively call the
unmodified comparison at an intermediate, larger syntax state.

In the mixed-limit case write
\begin{equation}
                  X=\KK Lc(A),\qquad Y=\KK Me(B),\quad M<L,
                                                               \label{u:display:8-1}
\end{equation}
where $L$ has a nonzero limit exponent. Recursively compare $X$ with $B$.
If $X\le B$, then $Y>B\ge X$. If $X>B$, the separation certificate
below, with lower profile $h_M$, gives
$X>h_M(B)\ge Y$, by Lemma~\ref{u:lem:direct}. There is no equality branch
in this mixed-limit case.

\subsection{The simultaneous limit-cut certificate}
Whenever one original side of a comparison is a bare limit-exponent
kernel $X=\KK Lc(A)$ and its result is $X>S$, the algorithm also
constructs, for every requested fixed lower profile $P$ below $L$, a
computable modulus $N_{X,S}(P)$ such that
\begin{equation}
                   X(x)>P(S(x))\quad(x\ge N_{X,S}(P)). \label{u:display:8-2}
\end{equation}
This concerns generated $P/R$ expressions $S$, not arbitrary integers
below $X$: the assertion for $S=X-1$ would be false, and subtraction
does not occur in the generated language.

Here is the certificate's complete case analysis, following the same
strictly smaller children as comparison.
\begin{enumerate}
\item If $S=ax+b$, use direct growth \eqref{u:display:5-3} for the fixed profile
$z\mapsto P(az+b)$ at the actual argument $A\ge x$. Positivity of $P$
makes this imply \eqref{u:display:8-2}.
\item If $S$ is nonlinear positive, $S<X$ forces every child peak $Z_i<X$
by the range-cell rule. Apply each smaller child's certificate to
$P\circ h_\omega$ and use $S<h_\omega(\max_i Z_i)$. Take the finite
maximum of their moduli; no permanent choice of maximizing child is needed.
\item If $S$ has the same outer principal $L$, strict order means either
a greater input for $X$, or equal inputs and smaller mass. The
lexicographic conclusion of Theorems~\ref{u:thm:adjacent} and
\ref{u:thm:limitmass} applies directly to $P$, at the genuine integer inputs.
\item If $S=\KK Me(B)$ has smaller outer principal, its comparison used
the smaller child $X>B$. Apply that child's certificate to
$P\circ h_M$ and use $S\le h_M(B)$. This handles any finite nesting of
lower contexts by recursive composition, not by an unproved absorption step.
\item If $S$ has larger limit outer principal $H>L$, its comparison with
$X$ compares $S$ against $A$. The only branch giving $X>S$ has $S\le A$;
direct growth for $\KK Lc$ gives $X>P(A)\ge P(S)$. In the other branch
the larger kernel's child certificate would force
$S>h_L(A)\ge X$, so cannot yield $X>S$.
\item A larger successor outer principal on $S$ is unfolded exactly until
it reaches $L$ or a larger limit exponent; use the corresponding preceding
case. The rewrites are exact, including their changed masses and domains.
\end{enumerate}
Every requested profile remains a fixed positive expression with heads
below the relevant limit $L$: only affine maps, fixed lower heads and
$h_\omega$ are composed. Thresholds are profile dependent, and are computed
by following these finitely many smaller children.

\subsection{A common well-founded measure}
Expand sharing when counting occurrences. Let $h$ be the number of
original source $h_\lambda$ nodes across the two compared expressions,
including their argument nodes. Let $\mathcal Q$ be the finite collection
of principal indices in the Cantor decompositions of all original source
heads and initial auxiliary kernels. Close it under finite successor-exponent
chains, stopping at exponent one or a limit exponent. No rule introduces
a principal outside this finite set. List it increasingly as
$\pi_0<\cdots<\pi_{m-1}$ and assign a kernel occurrence of index $\pi_i$
weight $\omega^{m-1-i}$. Let $W$ be their natural sum across the pair.
Use the lexicographic measure
\begin{equation}
                             (h,W).                 \label{u:display:8-3}
\end{equation}

Replacing an original tag by its peak removes that displayed $h$ node,
retaining its argument only once in each child comparison. In $P/P$ and
$R/P$ children, other siblings are not duplicated. Thus $h$ decreases.
Equal-principal comparison removes both outer kernels, decreasing $W$.
Mixed-limit comparison removes the lower outer kernel, also decreasing $W$.
A combined mixed-successor step removes the opposite lowest-rank $M$
occurrence and adds only finitely many occurrences of principals strictly
greater than $M$, retaining the original higher occurrence internally.
Their total weight is strictly smaller than one $M$ weight. Hence $W$
decreases. Affine or closed leaves stop.

Since $h$ is natural and $W<\omega^m$, every branch terminates. Branching
is finite, so the comparison tree is finite. The certificate uses that
same tree: each invocation in its case analysis belongs to a strictly
smaller comparison child or is a direct supplier application. Thus no
unchanged comparison state calls its own certificate. Mutual induction
on \eqref{u:display:8-3} proves both comparison correctness and \eqref{u:display:8-2}, establishing
Theorem~\ref{u:thm:sign}.

\subsection{Closure and exact finite exceptions}
The two sorts are closed under every recursive child. Peak replacement
introduces a finite kernel word on a positive argument plus its fixed
norm constant. Successor expansion introduces only a longer kernel word.
No signed inner offset, exposed count, division, inverse iterate or
variable ordinal index occurs. Increasing a mass changes fixed domain
constants without changing the finite rank set or measure.

\begin{example}[A nested comparison with the counts retained]
Consider $E=h_{2\omega}(x)$ and $F=h_\omega(h_\omega(x))$.
Deleting the two $\omega$ blocks gives
\[
 E=f(f(x+4))+2,\qquad F=f(f(x+2)+3)+1.
\]
Thus their peaks are respectively $f(f(x+4))$ and
$f(f(x+2)+3)$. The equal-outer-$\omega$ rule compares the arguments
$f(x+4)$ and $f(x+2)+3$. The latter positive nonlinear expression lies
strictly between its peak $f(x+2)$ and the next range value on the
computed tail. Comparing $f(x+4)$ with that peak reduces to the affine
comparison $x+4>x+2$. Hence $f(x+4)>f(x+2)+3$, and the original signed-row
rule returns $E>F$ on the effective tail. The offset $3$ has been paid
by a range-cell test, rather than deleted as negligible without proof.
The final procedure evaluates the finitely many earlier inputs exactly.
\end{example}

\begin{example}[Why the mass key changes at degree three]
At the same genuine input, the masses $1$ and $2$ give equal
$\omega^2$ kernels because their ceilings after division by two agree.
For a parent $\omega^3$ comparison with mass difference one, the first
common $\omega^2$ block can likewise tie when its retained child masses
are an odd $a$ and $a+1$. The next block reduces both child masses by
three; its ceiling classes are then distinct. The guaranteed two common
blocks ensure that those reduced masses remain nonnegative.
This is the two-block mechanism in
Theorem~\ref{u:thm:mass}, and explains why the algorithm uses neither strict
mass order at degree two nor a ceiling equivalence at all higher degrees.
\end{example}

\begin{corollary}[One-input quantified comparison]\label{u:cor:boolean}
For any finite Boolean combination of equality, weak-order and strict-order
atoms between expressions of $\mathcal P_1$, each of its existential and
universal closures over $x\in\NN$ is decidable.
\end{corollary}
\begin{proof}
Apply Theorem~\ref{u:thm:sign} to every atom and take the maximum tail bound.
The Boolean truth value is constant beyond that bound. Evaluate the
original positive input expressions exactly at every earlier positive
integer. For an existential statement, a true prefix instance or a true
tail suffices; for a universal statement require all prefix instances and
the tail. Evaluate original expressions, rather than auxiliary kernels
before their raised genuine domains.
\end{proof}


\section{Finite-degree kernels and exact counters}
\label{m:sec:polynomial-kernels}
The next class requires counters as well as terminal norms. A dominant
approximation would lose cancellations needed for multivariate comparison.
For a fixed nonzero limit put
\begin{equation}\label{m:eq:J}
 J_\lambda(t)=\psi\bigl(\lambda+t-N(\lambda)\bigr),
 \qquad t\ge N(\lambda),\qquad J_0(t)=t.
\end{equation}
Thus $h_\lambda(n)=J_\lambda(N(\lambda)+n)$, and $t$ is the
whole ordinal norm.  Corollary~\ref{f:cor:closed-evaluation} gives
$J_\lambda(t)=f^\tau(t)+\tau$, where $\tau$ is the complete
replacement count.  The offset $N(\lambda)$ is retained throughout.

For integer $d\ge1$, retained numerical norm $c\ge0$, and initial control
input $t$ with $c\le f(t)$, start with suffix $\omega^d$.
At step $i\ge1$ use total norm $f^i(t)$. If the nonfinite suffix is $\eta+n$, with nonzero limit
part $\eta$ and finite $n$, replace it by $M(\eta,f^i(t)-c)$: the greatest
suffix below its \emph{limit part} of norm at most $f^i(t)-c$.
This is the limit-part convention of the principal-kernel definition; taking
the predecessor below the whole current suffix defines a different process.
Stop as soon as the suffix is finite. Let $E^d_c(t)$ be the number of
steps and
\[
 X^d_c(t)=f^{E^d_c(t)}(t).
\]
Each replacement strictly decreases an ordinal, so the process terminates.
Theorem~\ref{f:thm:maximum-predecessor} supplies the
finite algorithm and full-budget saturation. The numerical definition makes
sense even for $c$ not realized by any actual ordinal prefix.
It also permits $c>t$: the initial $t$ records the control input, rather
than the norm $c+N(\omega^d)$ of the unpadded principal and numerical mass.
The first replacement has suffix norm $f(t)-c$ by limit-part saturation,
so its total norm is exactly $f(t)$; subsequent total norms are $f^i(t)$.
Writing $Q=M(\omega^d,f(t)-c)$, the complete count is one plus the
compensated controlled count starting at $Q$.

When $t\ge c+d+1$, the numerical process agrees, from its first replacement
through the first finite exit, with the principal kernel's norm-padded startup
$\omega^d+(t-c-N(\omega^d))$. This stronger condition is required for
correspondence with the principal kernels and the genuine-prefix interpretation; it is not
imposed on the broader numerical process.

For a genuine prefix $\mu$ whose least CNF exponent is at least $d$, and
$c=N(\mu)$, splitting the controlled history gives
\begin{equation}\label{m:eq:prefix}
 J_{\mu+\omega^d}(t)=J_\mu(X^d_c(t))+E^d_c(t),
       \qquad t\ge c+d+1.
\end{equation}
The output of $X$ is still the \emph{total} norm, so the prefix mass is not
subtracted again. An admissible signed pair $(d,j)$ has $d\ge2$ and
$j\in\mathbb Z$, or $d=1$ and $j\ge0$. For each such fixed pair, put
$X^d_{c,j}(t)=f^{E^d_c(t)+j}(t)$ on a tail where the exponent is positive.
For $d=1$, the depth is $1+j$, so a negative $j$ has no positive-depth tail.
These are coherent depths: $f(X^d_{c,j})=X^d_{c,j+1}$. A negative $j$ is a
fixed number of already available predecessors under $f$; it is not a new
variable-iteration graph.

\begin{lemma}[Exact block recursion]\label{m:lem:kernel-recursion}
One has
\[
 E^1_c=1,\quad X^1_c=f,\qquad
 E^2_c(t)=f(t)/2-\lceil c/2\rceil+1.
\]
Define
\begin{equation}\label{m:eq:V}
 V_j(t)=f^{\,f(t)/2+j+1}(t).
\end{equation}
Then $X^2_{c,j}=V_{j-\lceil c/2\rceil}$ for every fixed signed $j$ on
its positive-depth tail. For $d\ge2$, the unsigned identity
$X^d_c=X^{d-1}_c\circ X^d_{c+d}$ and the count identity
\eqref{m:eq:Erec} hold on a computable tail. The signed identity
\eqref{m:eq:Xrec} holds when $d\ge3$ or $j\ge0$:
\begin{align}
 X^d_{c,j}&=X^{d-1}_{c,j}\circ X^d_{c+d},\label{m:eq:Xrec}\\
 E^d_c&=E^{d-1}_c\circ X^d_{c+d}+E^d_{c+d}.\label{m:eq:Erec}
\end{align}
\end{lemma}
\begin{proof}
For $d=1$ the first predecessor is finite. For $d=2$, the first
predecessor has $\lfloor(f(t)-c)/2\rfloor$ copies of $\omega$ and a finite
remainder; each $\omega$ copy then costs one replacement. The function $f$
is even, which gives the displayed formula.
For general $d$ write $B=f(t)$ and $B-c=dq+\rho$, $0\le\rho<d$. The first
predecessor has leading part $\omega^{d-1}q$. Its remainder is zero if
$\rho=0$, finite one if $\rho=1$, and $\omega^{\rho-1}$ if $\rho\ge2$.
This is the largest remainder \emph{below $\omega^{d-1}$} under the norm
budget, not a maximum with no ordinal bound. Removing it leaves norm $B$
with spent count one in the first two cases, or norm
$X^{\rho-1}_{B-\rho}(B)$ with count $1+E^{\rho-1}_{B-\rho}(B)$ otherwise.
The remaining $q$ degree-$(d-1)$ blocks have retained masses
$c+d(k-1)$, for $k=q,q-1,\ldots,1$. Their $X$ maps compose and their $E$
counts add. Replacing $c$ by $c+d$ preserves the initial remainder and
initial entry but stops one block earlier. This proves both recurrences,
including the coherent final depth $j$ whenever its degree-$(d-1)$
child is admissible. For $d=2$ and negative $j$, the child has no positive
depth; the exact quadratic $V$ formula applies instead. This is also a
rank-recursive exact evaluator for every displayed kernel and counter.
\end{proof}

\begin{example}[The first counter that must be retained]
For $\lambda=\omega^3$, formula \eqref{m:eq:prefix} gives
$J_\lambda(t)=X^3_0(t)+E^3_0(t)$. For $\lambda=2\omega^3$, it gives
$J_\lambda(t)=J_{\omega^3}(X^3_4(t))+E^3_4(t)$, since
$N(\omega^3)=4$. Replacing either expression by its outer $X$ alone would
change exact graph equations and could change the result of cancellation.
\end{example}

\subsection{Separation with indices allowed to grow}
We need a uniform estimate, because the block recursion creates retained
masses comparable with $f(t)$. Merely proving a tail for each fixed mass
would not license the induction.

\begin{lemma}[Uniform kernel separation]\label{m:lem:uniform}
Put $C_d=1000d^2$, $T_3=16$, and for $d\ge4$ put
$T_d=\max(T_{d-1},100C_{d-1}+100d)$. For $d\ge3$, $c,e\ge0$, signed
$j,k$, and $R\ge0$, suppose
\begin{equation}\label{m:eq:region}
 t\ge T_d,\qquad
 f(t)\ge C_d(c+e+|j|+|k|+R+1).
\end{equation}
Then the depths are defined and at least $f(t)$, and
\begin{align*}
 X^d_{c,j}(t+1)&>f^R(X^d_{e,k}(t)),\\
 c<e\quad&\Longrightarrow\quad
 X^d_{c,j}(t)>f^R(X^d_{e,k}(t)).
\end{align*}
At equal input and equal $c$, exact order is order of $j$. On a common
computed tail, equality of two rank-$d$ values occurs exactly when input,
$c$, and $j$ are identical.
\end{lemma}
\begin{proof}
We first record the elementary degree-two estimate. If
$A=f(w)/2\ge |j|+|k|+R+4$, then
\[
 V_j(w+1)>f^R(V_k(w)).
\]
Indeed $f(w+1)/2-A\ge A$, so the left iterate starts at a larger seed and
has more iterations than the right one plus $R$. At equal input, the $V$
values are ordered exactly by their indices. This argument permits indices
and required depth varying with a later induction parameter, provided the
displayed numerical bound is paid.

For $d=3$, Lemma~\ref{m:lem:residues} at moduli two and three gives
that $B=f(t)$ is even and $B\equiv1\pmod3$; write
$a=(B-1)/3$, which is odd. The three initial entries for prefix classes
$\delta=0,1,2$ are
\[
\begin{array}{c|ccc}
\delta&\text{initial }c&\text{entry norm}&\text{spent count}\\ \hline
0&B-1=3a&B&1\\
1&B=3a+1&B&1\\
2&B-2=3(a-1)+2&f(B)&2.
\end{array}
\]
Descending $c$ by three applies $V_{-\lceil(c-3)/2\rceil}$ to the preceding
entry and adds its displayed iteration length to the counter. Write
$Y^\delta_r=X^3_{3r+\delta}(t)$. At the first common entry $r=a-1$,
$Y^0_r=f(Y^1_r)$, since the two subtracted ceilings differ by one. One
further aligned block amplifies this integer input gap to any fixed
$f$-depth by the degree-two estimate. Subsequent indices are favourable:
the class-zero index is no smaller than the class-one index. Next,
$Y^2_{a-1}=f(B)$ whereas
\[
 Y^1_{a-1}=f^\ell(B),\qquad
 \ell=f(B)/2-\lceil(3(a-1)+1)/2\rceil+1.
\]
For $a$ larger than the desired depth by ten, $\ell$ exceeds that depth
plus one. Finally align class two at entry $r-1$ with class zero at entry
$r$. The initial seeds satisfy $f(B)>B$, and
$\lceil(3r-4)/2\rceil\le\lceil(3r-3)/2\rceil$, so the degree-two
estimate again applies in the favourable direction. Consequently
\[
 Y^0_r>f^L(Y^1_r),\quad
 Y^1_r>f^L(Y^2_r),\quad
 Y^2_r>f^L(Y^0_{r+1})
\]
at all required target entries, after paying the fixed depth $L$.
All smaller block entries are at least $B$, and the absolute degree-two
indices are at most $\lceil3a/2\rceil+2$ plus fixed alignment shifts;
$f(B)/2$ pays their sum and $L$. Condition \eqref{m:eq:region}, with
$C_3=9000$, gives $c,e\le B/9000$ and
$L=R+|j|+|k|+1\le B/9000$.  Each target coefficient index is at
most $B/27000$, whereas $a=(B-1)/3$.  Thus there are more than
$L+10$ blocks between the starting classes and every required target.
The displayed quadratic estimate pays the signed depth offsets as well
as $R$.  This proves the same-input statement by chaining the three
adjacent classes.

For cross-input cubic comparison, $B'=f(t+1)>2^B$ supplies more than one
extra leading block. At the smaller process's starting coefficient,
the larger process has completed a quadratic block and has seed at least
$f(B')>f(B)$, exceeding either possible smaller seed. Align the remaining
blocks, shifting the coefficient alignment by at most
$\max(\lfloor c/3\rfloor-\lfloor e/3\rfloor,0)$. Their indices are bounded
by $\lceil3a/2\rceil+2$ plus this shift, and their differences are bounded
by twice the shift plus three. The same degree-two estimate at smaller
entry at least $B$ pays each comparison. Extra remaining blocks only
increase the final norm. This proves the cubic cross-input statement.

For the induction $d\ge4$ set $L=R+|j|+|k|+1$, $B=f(t)$.
Condition \eqref{m:eq:region} leaves at least $B/(2d)>L+2$ leading
blocks. Every aligned lower-rank mass is at most $2B$, and both its pair
sum and each self-pair sum, together with $L+1$, are at most $5B$.
All smaller aligned entry norms are at least $B$. Since
$B\ge T_d\ge100C_{d-1}$,
\[
 f(B)\ge2^B>5C_{d-1}B.
\]
Thus the induction hypothesis is valid at the \emph{smaller} entry norm
even though the new masses may be comparable with $B$.

At equal $t$ and $c<e$, align the two descents at the initial leading
coefficient of the $e$ process. If no extra leading block was needed,
the $c$ process's remainder is at least as large, with smaller retained
mass. If extra blocks were needed, its last extra full block exceeds the
other initial remainder. These facts order their aligned seeds weakly.
To justify this seed comparison explicitly, maximal predecessors are
monotone in both ordinal bound and norm budget. If entry norms $v\ge w$
and masses $a,b$ satisfy
$f^i(v)-a\ge f^i(w)-b$ at every step, the first process stays above the
second at simultaneous steps, so takes at least as many steps and has at
least as large a terminal norm. Decreasing mass at equal entry is therefore
favourable. A larger entry also pays an adverse mass difference whenever
$f(v)-f(w)\ge a-b$, since successive iterate gaps do not decrease.
At aligned block $k$ the lower masses are $c+d(k-1)<e+d(k-1)$.
Equal seeds use lower-rank same-input separation; distinct seeds use
lower-rank cross-input separation. The estimates just proved pay both,
including lower self-pairs. One such block gives depth $L$, and the rest
preserve it.

At $t+1$, its first budget $B'>2^B$ gives an extra full leading block
before alignment. Its entry is at least $B'$ and retained mass at most
$B'$. The smaller initial remainder is below that block, has entry $B$,
and retained mass at most $B$. The inequality
$f(B')-f(B)>B'$ pays even an adverse retained mass; hence the larger seed
is strict. At subsequent aligned blocks the mass pair is bounded by
$4B$; lower-rank cross-input separation, including its self-pair bounds,
applies as before. This proves the cross-input statement. Monotonicity
at fixed mass extends adjacent comparison to arbitrary larger inputs.

Finally the number of remaining blocks exceeds $|j|+|k|+1$, so signed
depths leave a positive iterate and output at least $f(t)$. Taking the
maximum of the cross-pair and the two self-pair thresholds excludes every
unequal-input or unequal-mass equality. The coherent identity settles the
remaining equal-input, equal-mass case.
\end{proof}

\subsection{Counter leaders, exact ties, and residues}
For $r\ge3$ define $A_r=t$ and successively
\[
 A_{h-1}=X^h_{c+h}(A_h)\quad(h=r,r-1,\ldots,4),\qquad
 u=X^3_{c+3}(A_3),\qquad H=f(u).
\]
We call $H$ the counter's \emph{leader}; $u$ is its leader preimage.

\begin{lemma}[Counter expansion and collision]\label{m:lem:counter}
On an effective tail,
\begin{equation}\label{m:eq:counter-expand}
 E^r_c(t)=\frac H2+1-\lceil c/2\rceil
                  +\sum_{h=3}^r E^h_{c+h}(A_h),
 \qquad \frac H2<E^r_c(t)<\frac{3H}{4}.
\end{equation}
Suppose counters of ranks $r\ge s\ge3$ have equal leaders. Then they
have the same prefix index $c$. If $r=s$ their whole values are equal;
if $r>s$ the smaller-rank input is the corresponding $A_s$, and
\begin{equation}\label{m:eq:collision}
 E^r_c(t)-E^s_c(A_s)=\sum_{h=s+1}^r E^h_{c+h}(A_h).
\end{equation}
Every counter introduced on the right has rank strictly greater than $s$.
For distinct leaders, counter order agrees with leader order. Equal
leaders are ordered by counter rank, with equality at equal rank.
\end{lemma}
\begin{proof}
Iterate \eqref{m:eq:Erec} until the cubic term; its bottom quadratic count
is $f(X^3_{c+3}(A_3))/2+1-\lceil c/2\rceil$. This proves the exact
identity. A descent with initial norm $w$ and $L$ steps satisfies
$L\le f^L(w)-w$, since each iterate increases its integer argument.
Thus each residual above rank three is less than its next entry norm,
hence less than $u$. The cubic residual plus the fixed correction is also
less than $u$, since it is at most $u-A_3+1-\lceil c/2\rceil$.
It is positive on a computed tail: the number of cubic blocks tends to
infinity. Therefore
\[
 H/2<E^r_c(t)<H/2+(r-2)u<3H/4
\]
once $f(u)>4(r-2)u$ and the finite correction bound is paid.

Equality of leaders first cancels $f$, then compares the two outer cubic
kernels $X^3_{c+3}$. Lemma~\ref{m:lem:uniform} forces equality of their
prefix indices and arguments. Cancel the equal rank-four kernels next,
then rank five, up to rank $s$. If $r=s$ the original inputs are equal.
If $r>s$ the rank-$s$ input is precisely $A_s$. Subtraction of the exact
recursions gives \eqref{m:eq:collision}. Its right side is positive and tends
to infinity. For different leaders $H_1<H_2$, their being distinct values
of $f$ gives $H_2\ge2^{H_1+1}$; the half-core bounds make the order strict.
This also proves the assertions with arbitrary fixed outer shifts after
raising the computed tail.
\end{proof}

\begin{lemma}[Effective residues of kernels and counters]\label{m:lem:kernel-residues}
For each fixed $d\ge1$, $c\ge0$ and $m\ge1$, the counter $E^d_c$
has an effectively computable constant residue modulo $m$ on a computable
tail. For every admissible signed pair $(d,j)$, the same holds for
$X^d_{c,j}$ on its positive-depth tail. The $X$
residue is $r_m$. A valid counter residue is
\begin{equation}\label{m:eq:E-residue}
 E^d_c(t)\equiv
 p_d\bigl((r_{Q_d}-c)\bmod Q_d\bigr)\pmod m,
 \qquad Q_d=d!m^{d-1},
\end{equation}
where $p_d(n)$ counts partitions of $n$ into parts at most $d$.
\end{lemma}
\begin{proof}
A coherent $X$ lies in the range of $f$ with preimage at least $t$ on a
computed tail. Apply Lemma~\ref{m:lem:residues}.
For the counter we use the restricted-partition and congruent-budget
theorems of \cite{wwcounts}, with the following exact contract.  The
integer $Q_d=d!m^{d-1}$ is a period of $p_d(n)$ modulo $m$.  A
norm-padded $\omega^d$ block of initial total norm $t\ge c+d+1$,
retaining numerical mass $c$, has complete count congruent to
$p_d(S-c)$ modulo $m$ if every subsequent total budget is at least
the current total norm and congruent to one $S\ge c$ modulo $Q_d$.
The predecessor is taken below the current limit part and the segment
ends at the first finite suffix.  The budgets need not be equal.

Take
\[
 t\ge\max\{c+d+1,Q_d\},\qquad
 S=c+((r_{Q_d}-c)\bmod Q_d).
\]
The numerical kernel agrees with the genuine norm-padded block on this
tail.  Every actual budget $f^i(t)$, $i\ge1$, exceeds the preceding
total norm and is congruent to $r_{Q_d}$, by
Lemma~\ref{m:lem:residues}.  The cited transfer theorem therefore applies.
The partition period gives \eqref{m:eq:E-residue}.  Finite dynamic
programming computes $p_d(n)$ for $n<Q_d$; all smaller valid inputs
can be evaluated exactly.  This proves an explicit residue tail without
using comparison or an existential decision procedure.
\end{proof}

\section{Closed existential word arithmetic}
\label{m:sec:word-arithmetic}
Fix a finite maximum degree $D\ge3$. A \emph{word} in one raw variable
is a finite composition of signed constant shifts and the functions
\[
 f,\quad V_j,\quad X^r_{c,j},\quad E^r_c\quad(3\le r\le D),
\]
with fixed integer indices satisfying the preceding conventions. A
word-linear row is an integer linear combination of finitely many such
words, evaluated in possibly different raw variables, plus a constant.
We allow arbitrary Boolean combinations of row comparisons with zero and
fixed-modulus congruences. Each raw variable has an explicit individual
lower bound ensuring that all its words are defined.
These bounds concern every displayed occurrence, including an occurrence
whose coefficient is zero.  The domain of the whole query is the
conjunction $\bigwedge_i x_i\ge b_i$ of those individual bounds;
it is not changed by a Boolean negation inside the matrix.

Every word has a computable domain tail and a lower offset $M$ with
$W(n)\ge n-M$. This follows structurally: shifts add to the offset;
$f$, $V$, and coherent $X$ dominate their input on their stated tails;
counters do so as well. For the latter, a sufficiently large degree-three
block contains two quadratic blocks, and a quadratic block entered at
least at $f(t)$ already spends more than $f(t)$ steps. Higher-degree blocks
contain one of these. These structural bounds also pay every threshold at
an interior argument of a composition. They are not obtained by treating
the interior argument as a raw variable.

\begin{lemma}[Ordered maximum-variable guards]\label{m:lem:guards}
For two words $A(x),B(y)$ and the guard $x\ge y$, an effective common
tail permits their three-way comparison to be rewritten as a finite Boolean
combination in which $x$ occurs only linearly and words are applied only
to the remaining variable $y$. For two words at the \emph{same} raw
variable the corresponding tail comparison is effectively decided.
All thresholds and newly generated words are computed before the tail is
used.
\end{lemma}
\begin{proof}
Peel fixed outer shifts. Unequal range values have gaps eventually larger
than a fixed shift. Lemma~\ref{m:lem:counter} gives the same fact for two
counter values, including equal-leader unequal-rank values whose exact
residual difference tends to infinity. For a counter with preimage $u$
and range value $f(v)$, its half-core bounds give
\[
 E^r_c<f(v)\ \Longleftrightarrow\ u\le v,
 \qquad E^r_c>f(v)\ \Longleftrightarrow\ u>v;
\]
the gaps again tend to infinity. Thus fixed shifts are kept for exact-core
ties and otherwise cannot close the gap.

For $f/f$ strip the two $f$ nodes. For $f$ against a coherent range
kernel, move one fixed depth to its index and strip $f$; for example
$V_j=f\circ V_{j-1}$. Two $V$ heads compare lexicographically by argument,
then by index, using the elementary estimate in the proof of
Lemma~\ref{m:lem:uniform}. Two equal-rank $X$ heads compare by argument,
then by $(-c,j)$, using that lemma and retaining the exact equal-input
branch.

For $X^r$ against $X^s$ with $r>s\ge3$, expand only the higher-rank head
by \eqref{m:eq:Xrec} until its outer rank is $s$, and immediately compare
and strip that head against the opposite rank-$s$ head. If $s=2$,
the exposed quadratic head is $V_{j-\lceil c/2\rceil}$ and the same
procedure uses its index comparison. No inverse counter is introduced.
For counter/range comparison use the displayed preimage rule. For
counter/counter comparison use their leader preimages followed by the
rank and exact-tie rule of Lemma~\ref{m:lem:counter}.

Here is a termination measure, needed because mixed-rank expansion can
increase the written expression. First count all $E$ nodes in the two
words. Replacing a counter by its leader-preimage composition removes
its outer $E$ node and retains its argument once. For the range part,
assign each rank-$s$ occurrence weight $\omega^{D-s}$, where $f$ has
rank one and $V$ rank two, and naturally sum these weights. A mixed
comparison deletes the opposite rank-$s$ occurrence and introduces only
finitely many occurrences of ranks strictly above $s$, retaining the
original higher-rank occurrence. Its ordinal weight decreases. Equal-rank
stripping and $f$ cancellation decrease weight as well. The lexicographic
pair (number of $E$ nodes, range weight) is therefore well-founded along
every recursive child. There are finitely many branches at each step.

Ultimately the $x$ side is a shifted leaf, so its comparison with a word
at $y$ is linear in $x$; or the $y$ side is a shifted leaf and the $x$
side is nonlinear. In the latter case the nonlinear value is at least
$2^{x-M}$ for a computed $M$, so eventually exceeds $y$ plus any fixed
shift, since $y\le x$. When both words have the same raw variable,
continue stripping until leaves or leaf/nonlinear growth decides the
comparison; do not create a purported linear atom $x<W(x)$.
Take the finite maximum of all local bounds using the structural offsets
of interior arguments. The finite recursion tree, rather than a guessed
fixed envelope of indices, computes that maximum.
\end{proof}

The guards compare individual word values. A linear row can still cancel
its largest values, so it requires a different induction.

\begin{lemma}[Ranked cancellation]\label{m:lem:row}
For a finite collection of word-linear rows and a chosen maximal raw
variable $x$, one can compute a tail and finite guarded cases in which
each row either has a known strict sign or is exactly linear in $x$.
The guards can themselves be made linear in $x$ by
Lemma~\ref{m:lem:guards}. No correlation between word values is discarded.
\end{lemma}
\begin{proof}
Absorb fixed outer shifts into the constant and keep raw-linear terms
separate. The other top-level atoms are ordinary range values (rank zero
for this row argument) and counters of ranks $3,\ldots,D$. Partition their
ordinary values and counter leaders into ordered equality classes, retaining
every equality and order guard. Merge equal-rank counters in a class as
\emph{whole counters}, justified by Lemma~\ref{m:lem:counter}. Prefer a
representative actually involving $x$ when one remains in the current
class. A class now has ordinary coefficient $a$ and at most one coefficient
$b_r$ for each counter rank. Its leading coefficient is
\[
 a+\tfrac12\sum_r b_r.
\]
If it vanishes, pay its residuals before deciding any sign. When $a\ne0$,
expand the counters by \eqref{m:eq:counter-expand}; the common half-leader
cancels the ordinary atom, leaving only residual counters and constants.
When $a=0$, let $s$ be the least surviving counter rank. Rewrite each
higher counter by \eqref{m:eq:collision} as the common rank-$s$ counter plus
higher-rank residuals. Since $\sum b_r=0$, the common term cancels. Remove
an entirely zero class. Repartition after every rewrite, since its newly
exposed leaders may coincide with leaders elsewhere in the row.

Count \emph{top-level} row atoms, assigning the ordinary rank-zero atom
weight $\omega^D$ and a rank-$r$ counter weight $\omega^{D-r}$.
Natural sum is a well-founded measure. The first rewrite loses a rank-zero
atom and adds only positive-rank atoms. The second loses the least rank
$s$ and adds only ranks strictly greater than $s$. Existing higher-rank
atoms can proliferate, but cannot offset that ordinal decrease. Nested
nodes inside their arguments are not included in this measure; their
comparisons terminate by the separate guard measure. Every branching row
rewrite therefore strictly decreases weight. The tree is finitely
branching and has no infinite branch, so it is finite. Enumerating it to
completion computes all generated words, coefficients, prefix indices,
and guard thresholds before any finite-prefix case is split off.

After the rewrites every surviving leading coefficient has absolute value
at least $1/2$. Let $P=f(u)$ be the greatest surviving leader, $S$ the
current absolute nonlinear coefficient mass, and $A,B$ bounds for the
absolute raw-linear coefficients and constant. The exact counter
corrections have total size at most $DSu$. Every smaller leader $Q$
satisfies $Q<\log_2 P$, so its total contribution is at most
$S\log_2P$. For $u\ge8DS+16$ each of these two errors is less than $P/16$.

If a current nonlinear $x$ occurrence survives, the greatest leader is
at least its leader. Structural offsets therefore give $u\ge x-M$ for a
computed $M$. Require $x\ge2M+16$ and $x\ge8(A+B+1)$. The elementary
inequality $2^{x-M}\ge x^2$ on this range, proved at the initial value
and then by doubling, gives $Ax+B<P/8$. Total error is less than $P/4$,
while the leading magnitude is at least $P/2$. Hence its coefficient
strictly determines the sign. If no current nonlinear $x$ occurrence
remains, keep the exact row, which is linear in $x$.

Growth provenance refers only to current occurrences. A deleted leader
passes no $x$ label to a residual which does not actually contain $x$.
This prevents an invalid sign decision after cancellation has removed all
nonlinear dependence on $x$. Take maxima of these explicit bounds over the
finite row and guard trees. Ordered partitions retain the original shared
arguments throughout, even when an inconsistent branch is later discarded.
\end{proof}

\begin{theorem}[Closed existential word arithmetic]\label{m:thm:existential}
The existential closure of every Boolean combination of word-linear row
comparisons and fixed-modulus congruences in the finite-degree language
above is decidable, on its explicit word domains.
\end{theorem}
\begin{proof}
First remove a variable which occurs only linearly, if any, by ordinary
Presburger quantifier elimination \cite{cooper}, including its explicit
lower bound in the matrix. Words in other variables are parameters
in this step; the eliminated variable occurs in none of them.  The
Presburger equivalence is valid for every assignment of its parameter
symbols, hence for their actual correlated word values.  Its output
uses only integer-linear rows and fixed-modulus congruences in those
same parameters.  Substituting the words back therefore stays in the
stated syntax and does not apply any word to a new sum. Repeat or
evaluate the closed matrix if no variables remain.

Otherwise branch on a maximal raw variable $x$, retaining $x\ge x_i$.
Compute the complete finite guard and row trees of
Lemmas~\ref{m:lem:guards}--\ref{m:lem:row}. Compute residue thresholds for each
congruence as well: Lemma~\ref{m:lem:kernel-residues}, composed with interior
lower-offset bounds, makes every nonlinear word's residue constant.
Collect shifts in its fixed residue. On that tail, congruences are thus
congruences in raw variables only. Take a common $N$ paying all original
and generated domains, sign estimates, guards, and residue bounds.

Before executing any tail transformation, split the \emph{original current
matrix} into the all-raw-inputs-$\ge N$ branch and finitely many branches
fixing some raw $x_i=c<N$ above its original lower bound. In a finite
branch substitute $c$ into its original words and evaluate them exactly.
Do not evaluate a newly introduced coherent negative-depth word below the
tail on which it was introduced. Exact kernel evaluation is finite by the
block recursion. Each such branch removes a quantified variable.

On the all-large branch, use the guarded row and congruence transformations.
Everything now contains $x$ only linearly. Presburger elimination of
$\exists x$, treating words in the remaining variables as parameters,
again produces word-linear rows and congruences with one fewer variable.
Preserve all Boolean connectives and every equality/order guard. No word
has been applied to a sum, and no auxiliary derived value has been made
an independent witness.

The output retains the lower bounds of the remaining variables.  Each
new word introduced in the finite normalization tree has already had
its genuine domain paid by the chosen common $N$.  Thus the recursive
call has the same explicit individual-domain contract as the input.
Any new congruence modulus created by Presburger elimination is dealt
with in that smaller-variable call, not by restarting the present
dimension in search of a stable cutoff.

The induction measure for the whole algorithm is the number of raw
existential variables. It strictly decreases in every finite-prefix,
pure-linear, and all-large branch. The subsidiary guard and row trees are
finite by their own measures. Consequently this is a terminating effective
procedure. No feasible complexity estimate follows: finite exception values
and arithmetic thresholds can be extremely large.
\end{proof}

This theorem is deliberately about \emph{closed existential} sentences.
In a formula with a free parameter, that parameter could be the largest raw
argument and could not be removed by the stated induction. We have proved
neither free-parameter quantifier elimination nor the full first-order
theory of arbitrary polynomial-collapse functions.

\subsection{Exact polynomial graphs}
\begin{theorem}[Polynomial collapse graphs]\label{m:thm:graphs}
For each fixed nonzero polynomial limit $\lambda<\omega^\omega$, the graph
$y=h_\lambda(n)$ is effectively expressible as a finite disjunction of
word-linear equations in $n,y$, with explicit domains and exact finite
exceptions. One finite degree bound suffices for each finite collection of
such graphs. Arbitrary fixed nesting and addition of these functions admit
a total functional existential graph translation.
\end{theorem}
\begin{proof}
Remove one unit at the least positive degree of $\lambda$, writing
$\lambda=\mu+\omega^d$, and put $c=N(\mu)$. There are finitely many
monomial units. Starting with $2J_0(t)=2t$, compile \eqref{m:eq:prefix}
recursively: substitute the unary word $X^d_c(t)$ for the input in the
already compiled expression for $2J_\mu$, and add $2E^d_c(t)$. For $d=1$
use $X=f$, $2E=2$; for $d=2$ use
\[
 X=V_{-\lceil c/2\rceil},\qquad
 2E=f(t)-2\lceil c/2\rceil+2;
\]
for $d\ge3$ retain the $X$ and $E$ heads. This substitution composes unary
words only. It never applies a function to a sum of compiled values.

Finally replace $t$ by $n+N(\lambda)$. On a computed tail this yields
$2y=P_\lambda(n)$ with an integer word-linear $P_\lambda$. The finitely
many excluded $n\ge0$ are evaluated by the total closed predecessor
algorithm and included as disjuncts $n=k$, $y=h_\lambda(k)$. The graph is
therefore total and functional on the original natural domain.
To flatten a nested expression, introduce a fresh natural variable for
each arithmetic node: addition and fixed multiples have linear equations;
a collapse node has the graph just constructed, with its original argument
variable. Preserve all repeated original variable names and quantify every
new node variable. This creates no independence assumption: the graph
equations enforce precisely the original correlations.
\end{proof}

\begin{lemma}[From finite graph exceptions to word domains]
\label{m:lem:graph-domains}
Let finitely many total functional graph clauses be obtained from
Theorem~\ref{m:thm:graphs}, together with linear arithmetic graph
equations and a Boolean combination of linear comparisons between graph
wires.  Their closed
existential conjunction is decidable by
Theorem~\ref{m:thm:existential}, with original positive inputs and
fresh natural graph variables.
\end{lemma}
\begin{proof}
Every collapse graph is a finite disjunction of its tail equation,
guarded by a computed input lower bound, and finitely many exact
small-input point clauses.  Distribute only these finitely many
positive graph disjunctions.  In a resulting case, a small-input
clause fixes its input and output to known integers.  Substitute such
equalities into the case.  If a retained tail guard on the same input
fails at the substituted integer, discard the inconsistent case;
otherwise any word in that fixed input is defined and can be evaluated
exactly.  Repeat the finite substitutions if desired.

All remaining word occurrences come from tail clauses with their
explicit bounds.  Take, for each shared raw variable, the maximum of
the bounds required by its remaining words, together with its original
natural or positive domain.  The accompanying guards assert exactly
these bounds.  The case is now a closed existential Boolean combination
of word-linear rows and congruences on the individual word domains, so
Theorem~\ref{m:thm:existential} applies.  The disjunction is finite.

This argument does not replace a correlated node value by a free
parameter.  All linear graph equations and original names remain in
each case.  It also does not negate a graph: the output test may contain
negations, but the functional graph constraints remain positively
conjoined.  Their uniqueness on the original domains gives the intended
source semantics.
\end{proof}

\begin{corollary}[Polynomial-head Boolean comparison]
\label{m:cor:polynomial-booleans}
For positive expressions in any finite number of raw inputs, using only
fixed polynomial limit heads below $\omega^\omega$, the existential
and universal closures of every finite Boolean combination of equality
and order atoms are decidable.  Arbitrary fixed nesting and addition
are permitted.
\end{corollary}
\begin{proof}
Flatten all the finitely many input circuits together, assigning one
fresh natural wire to each node and one shared input to each repeated
raw name.  Conjoin their total functional graphs positively and append
the Boolean output test.  The graph constraints have a unique extension
at every original positive valuation, so existence of this full
conjunction is equivalent to the original existential query.
Lemma~\ref{m:lem:graph-domains} decides it.  For a universal query,
negate only the Boolean output test and decide the resulting existential
counterexample; the graph conjunction remains positive.  All fresh
wires and original inputs belong to the same existential block.
\end{proof}


\section{Parameter-preserving fragments with arbitrary heads}
\label{m:sec:certificates}
Polynomial graphs were needed for completely unrestricted arrangements of
polynomial heads. Different arrangements permit arbitrary fixed heads
without constructing their graphs. This section isolates the spacing fact
used by all of them.

\begin{lemma}[Sparse endpoint comparison]\label{m:lem:endpoints}
For a fixed nonzero limit $\lambda$, $J_\lambda$ is strictly increasing.
Its distinct norm inputs have distinct terminal endpoints, ordered in the
same way. For any two distinct endpoints $Z<W$ one has
\begin{equation}\label{m:eq:sparse}
 W\ge2^{Z+1}.
\end{equation}
Every value with endpoint $Z$ lies strictly between $Z$ and $2Z$ and below
the next value of $f$.
\end{lemma}
\begin{proof}
Strict increase is controlled strictness from Lemma~\ref{f:lem:collapse-estimates}: when $t<u$, the smaller ordinal has norm
$t\le f(u)$, so the controlled comparison premise holds. Write
$J_\lambda(t)=f^L(t)+L$ and $J_\lambda(u)=f^M(u)+M$.
If their endpoints were equal, injectivity of $f$ would force $L>M$;
otherwise cancelling $L$ iterates would give $t\ge u$. The equal endpoints
and $L>M$ would imply $J_\lambda(t)>J_\lambda(u)$, a contradiction.
Put $s=f^{L-1}(t)$. Since $t\ge N(\lambda)\ge2$, one has $L\le s-1$,
so $f(s)<J_\lambda(t)<2f(s)<f(s+1)$ by Lemma~\ref{m:lem:growth}.
These intervals are disjoint and preserve order. If $W>Z=f(s)$ is another
endpoint, then $W\ge f(s+1)\ge2^{Z+1}$.
\end{proof}

\subsection{Nested common-head sums}
Fix any finite list of fixed nonzero limits $\lambda_1,\ldots,\lambda_l$.
Let $\mathcal A_0$ consist of natural constants, positive raw variables,
addition, and $f$, with arbitrary fixed nesting. Recursively set
\begin{equation}\label{m:eq:common-class}
 \mathcal A_r=\left\{c+\sum_i a_i h_{\lambda_r}(B_i):
 c\ge0,\ a_i>0,\ B_i\in\mathcal A_{r-1}\right\}.
\end{equation}
The sum can be empty. Each input supplies this finite level certificate.
There is no open remainder outside the head sum at a positive level.
For example $x+h_\lambda(y)$ does not pass that positive-level certificate
merely because both $x,y$ are allowed at level zero. Heads at different
levels may be completely different, and may exceed $\omega^\omega$.

\begin{theorem}[Common-head layers]\label{m:thm:common}
For $E,F\in\mathcal A_r$, each pointwise relation $E<F$, $E=F$, $E>F$
has an effectively constructed formula in $(\N;0,1,+,<,f)$ with the same
raw free variables and exactly the same truth at every positive valuation.
Hence arbitrary Boolean and quantified combinations of certified comparisons
are decidable.
\end{theorem}
\begin{proof}
Subtract the two expressions and write the row as
\[
 Q=K+\sum_{i<m}a_i J_{\lambda_r}(t_i),\qquad
 t_i=N(\lambda_r)+B_i.
\]
Coefficients and $K$ are signed integers. With no remaining atoms the sign
is the computable sign of $K$. Otherwise set
$C=\sum_i|a_i|$ and $L=\max(4,2C,4|K|+4)$.
On the branch with all $t_i\ge L$, partition the inputs by equal values,
combine coefficients, and delete zero blocks. If every block cancels,
$Q=K$. Otherwise its sign is the coefficient of the greatest surviving
input. Indeed let $W$ be its endpoint and $Z$ the greatest lower endpoint,
if any. Lemma~\ref{m:lem:endpoints} gives
$W\ge2^{Z+1}\ge2Z^2\ge4CZ$, since $Z\ge L\ge\max(4,2C)$.
All lower contributions have total absolute value less than
$2CZ\le W/2$. Also $|K|<W/4$, whereas the leading contribution has
absolute value greater than $W$. Its sign cannot change. Larger canceled
blocks play no part in this estimate. Enumerating the finitely many total
preorders of the labeled inputs produces a finite formula, guarded by
$t_i=t_j$, $t_i<t_j$, and $t_i\ge L$.

For the other branch choose a small remaining index, enumerate
$N(\lambda_r)\le c<L$, retain its exact guard $t_i=c$, evaluate the
closed $J_{\lambda_r}(c)$, fold $a_iJ_{\lambda_r}(c)$ into $K$, and
remove that occurrence. Recompute $L$ and repeat. Each small call removes
one occurrence; its possibly enormous updated constant is still a fixed
computable integer. Overlapping branches are harmless in a disjunction.
No open value is folded into a constant before its whole argument is fixed.

Induct now on $r$. At level zero comparisons are base-language atoms.
At higher levels every guard compares preceding-level expressions with
each other or with constants; a proposed constant below $N(\lambda_r)$
gives a constant truth value. Thus every generated predicate is covered
by the induction hypothesis. The lexicographic pair (level, remaining
head occurrences) proves termination: small folding decreases the second,
guard translation the first. Shared variables remain shared throughout.
The result is an exact base formula, to which
Theorem~\ref{m:thm:snyatkov} applies after any desired raw quantifiers are
added.
If cancellation removed a free name syntactically, conjoin the tautology
$x=x$ for that original name.  This gives exactly the original free-name
set without changing the predicate.
\end{proof}

For quantified combinations, translate the comparison atoms in place and
preserve the original Boolean structure.  A positive existential binder
becomes $\exists x\,(x\ge1\wedge\varphi)$, and a positive universal
binder becomes $\forall x\,(x\ge1\mathbin{\to}\varphi)$.  Bound names
can first be renamed to avoid capture.  Each atom may have its own
genuine finite certificate; its translation creates no independent raw
name for a derived value.  This is a parameter-preserving reduction to
the full base theory, so alternating binders are allowed here.

\begin{example}[Cancellation changes the controlling input]
Consider $2h_\lambda(x)+h_\lambda(y)-3h_\lambda(z)$. On the computed
large-input branch, a unique largest argument determines the sign by its
coefficient. But on the guard $x=z$, the row becomes exactly
$h_\lambda(y)-h_\lambda(x)$, whose sign is that of $y-x$; when all three
arguments agree the row is zero. The algorithm retains those guards rather
than treating three head values as independent large numbers. Small inputs
are evaluated by the finite folding procedure, so this is an exact
comparison theorem rather than only an eventual rule.
\end{example}

\subsection{Different heads in one control interval}
\begin{lemma}[Head order in a strict interval]\label{m:lem:head-window}
Let $\lambda>\mu$ be nonzero fixed limits. If
$t\ge\max(N(\lambda),N(\mu)+2)$ and
$N(\mu)\le u<f(t)$, then $J_\lambda(t)>J_\mu(u)$ and their endpoints
differ. Consequently, if finitely many norm inputs $t_i$ have minimum
$s\ge\max_iN(\lambda_i)+2$ and $\max_i t_i<f(s)$, their outputs are
ordered lexicographically by (ordinal head, norm input). Equality requires
equality of both, and different output blocks have distinct sparse endpoints.
\end{lemma}
\begin{proof}
The least limit above $\mu$ is $\mu+\omega$, of norm $N(\mu)+2$.
Controlled strictness at common norm $t$, or equality of the two heads,
gives
\[
 J_\lambda(t)\ge J_{\mu+\omega}(t)
       =J_\mu(f(t))+1>J_\mu(u).
\]
The exact middle identity follows by removing the final $\omega$ block:
its first predecessor retains $\mu$, assigns the remaining full budget
to its finite tail, and has total norm $f(t)$. If the two displayed
outputs had equal endpoints, their complete counts would satisfy
$L>M$; canceling the common $f$ iterates would give
$u=f^{L-M}(t)\ge f(t)$, a contradiction. For the finite collection, the
lower input is below $f(s)\le f(t)$ for every higher-head input $t$.
Apply this argument to different heads and
Lemma~\ref{m:lem:endpoints} within the same head.
\end{proof}

The interval inequality is strict for a reason. At $u=f(t)$,
\[
 J_{2\omega}(t)=f^2(t)+2,\qquad
 J_\omega(u)=f^2(t)+1.
\]
Their endpoints coincide, and
$J_{2\omega}(t)-2J_\omega(u)=-f^2(t)<0$ despite the positive coefficient
of the greatest ordinal head. An algorithm using head order at that
boundary would therefore be incorrect.

\subsection{Balanced heterogeneous layers}
The next certificate permits different sibling heads and an open lower-level
remainder. At level zero use the same $\mathcal A_0$ as before.  Define a
certificate recursively.  At a positive level it consists of one
preceding certificate $C$ and a finite ordered basis
$B_1,\ldots,B_k$, all certified by that same $C$.  The basis may have
repeated or dependent entries.  Expressions with the new certificate
are nonnegative affine combinations of the $B_j$ and atoms
\begin{equation}\label{m:eq:balanced-atoms}
 h_\lambda\!\left(c+\sum_j a_jB_j\right),
       \qquad c\ge0,\quad a_j\ge1.
\end{equation}
Each atom may have its own arbitrary fixed head and coefficients. The
essential condition is that \emph{every head argument includes every
common basis entry positively}. The basis is checked on the whole pair
being compared, including its outside remainder. A variable appearing
only in that remainder is not admitted by ignoring it while choosing the
basis. Zero multiples are discarded when computing actual supports.
For each fixed certificate its expression family is closed under
nonnegative affine combinations: collect constants, outside coefficients,
and atom lists.  Induction gives the same closure at the preceding
certificate.  Consequently the lower affine guards used below really do
belong to one preceding comparison problem; separate certificates are
not silently combined into an independent basis.

\begin{theorem}[Balanced comparison]\label{m:thm:balanced}
For two expressions with one finite balanced certificate, their three-way
pointwise comparison has an effectively constructed base formula in
$(\N;0,1,+,<,f)$ with exactly the same free variables. Arbitrary Boolean
and quantified combinations of such comparisons are decidable.
\end{theorem}
\begin{proof}
A signed comparison row has the form
\begin{equation}\label{m:eq:balanced-row}
 Q=K+\sum_jr_jB_j+\sum_{i<m}q_iJ_{\lambda_i}(t_i),\qquad
 t_i=b_i+\sum_j a_{ij}B_j,\quad b_i=N(\lambda_i)+c_i.
\end{equation}
The head-free case is already a signed affine comparison at the preceding
level. Otherwise put
\[
 S=\sum_jB_j,\quad C=\sum_i|q_i|,\quad C_0=\max_j|r_j|,
 \quad A=\max_{i,j}a_{ij},\quad B=\max_i b_i,
\]
with $C_0=0$ for an absent remainder and $A=1$ for an empty basis. Set
\begin{equation}\label{m:eq:balanced-L}
 L=\max\{16,2C,4C_0+4|K|+4,A+B+2,\max_iN(\lambda_i)+2\}.
\end{equation}
On the branch $s=\min_i t_i\ge L$ one has
\[
 S\le t_i\le AS+B,\qquad
 \max_i t_i\le As+B<s^2\le2^s\le f(s).
\]
Thus Lemma~\ref{m:lem:head-window} orders all head outputs by head and then
input. Partition the inputs within each equal head, combine coefficients,
and discard zero blocks. If all cancel, the remaining row is exactly the
preceding-level $K+\sum r_jB_j$.

If a block survives, let $W$ be the endpoint of the last one and $v$ its
norm input. Its coefficient has absolute value at least one, so its
contribution has magnitude greater than $W$. The sparse-endpoint estimate
bounds all lower head contributions by $W/2$, exactly as in the proof of
Theorem~\ref{m:thm:common}. The outside remainder is bounded using this
\emph{surviving} input, not a possibly canceled maximum:
\[
 \left|K+\sum_jr_jB_j\right|\le |K|+C_0v
          <v^2/4\le f(v)/4\le W/4.
\]
The surviving last coefficient therefore determines the strict sign.
The finitely many guarded preorders within equal heads express this rule
by comparisons of affine preceding-level expressions.

We must retain every small exception. Every positive expression in the
certificate satisfies $B(\vec x)\ge x$ for each raw variable in its
actual support. Induct through addition and $f$, and at a head use
$h_\lambda(e)>e$: the finite ordinal $e$ is below $\lambda+e$ and has
norm $e\le f(N(\lambda)+e)$, so controlled strictness and $\psi(e)=e$
apply. Fixed positive multiples preserve the property. If some $t_i<L$,
positivity of every $a_{ij}$ implies $S<L$. Therefore every raw variable
anywhere in the basis, including the outside remainder, is less than $L$.
Enumerate the finite box $[1,L-1]^V$ on its \emph{actual} support $V$;
evaluate each closed row and its inputs exactly, retaining precisely
points where some $t_i<L$ and the requested sign holds. Variables outside
$V$ remain unrestricted; empty support gives a single empty valuation.

All large-branch guards, and the all-canceled remainder, are affine
comparisons of the $B_j$ and thus have their preceding-level certificate.
The finite-point exception formula is already in the base language.
Induction on certificate level proves the claimed translation and
termination. Theorem~\ref{m:thm:snyatkov} supplies its arithmetic decision.
Repeated basis entries and shared variables never become independent
parameters in this construction.
Conjoining $x=x$ for any canceled original name makes the free-name
set exact.  The in-place Boolean and positive-binder translation following
Theorem~\ref{m:thm:common} then gives arbitrary quantified decisions.
\end{proof}

\begin{example}[A lower-head argument may be larger]
Let $E,F$ have a common preceding-level certificate and $\lambda>\mu$.
The row
\[
 3h_\lambda(E+F+2)-2h_\mu(2E+F+1)-F
\]
is balanced on basis $(E,F)$. Its two norm arguments need not have the
same order as their ordinal heads, and $F$ remains open outside them.
Nevertheless, once both norms exceed \eqref{m:eq:balanced-L}, the row is
positive by the higher-head coefficient. Its exact comparison is completed
by the finite support box. Omitting $F$ from one head argument would
remove the inference that a small input bounds the entire open remainder.
\end{example}

\subsection{One arbitrary core: direct Presburger comparison}
An especially useful specialization no longer requires a certificate for the
interior core. Let $E$ be any expression made from natural constants,
positive raw variables, addition, fixed nonnegative multiples, and arbitrary
fixed nonzero-limit functions $h_\lambda$, with arbitrary finite nesting.
No subtraction or exposed $f$ constructor occurs inside this stated grammar.
The only property needed below is $E\ge x$ for every variable in its actual
support, proved by the positive-source induction just given. It need not
have a balanced-level certificate, nor range independently over all naturals.

\begin{theorem}[Arbitrary single-core rows]\label{m:thm:single-core}
For the signed row
\begin{equation}\label{m:eq:single-core}
 Q(\vec x)=K+rE(\vec x)+
       \sum_i q_i h_{\lambda_i}(a_iE(\vec x)+c_i),
 \quad a_i\ge1,\ c_i\ge0,
\end{equation}
with fixed integer $K,r,q_i$ and fixed nonzero limit heads, each sign
predicate $Q<0$, $Q=0$, $Q>0$ has an effectively computable finite
Presburger formula in the original raw variables. Thus universal comparison
of these rows needs no arithmetic decision theorem about $f$.
\end{theorem}
\begin{proof}
Collect identical triples $(\lambda_i,a_i,c_i)$ and discard their zero
coefficients. If any heads remain, use \eqref{m:eq:balanced-L} with the
single formal basis $E$, $C_0=|r|$, and $b_i=N(\lambda_i)+c_i$. Increase
that bound to
\[
 T=\max\bigl(L,1+\max_{i,j}|b_i-b_j|\bigr).
\]
For $E\ge T$, the norms lie in one strict $f$ interval. Within equal
heads their affine inputs are ordered lexicographically by $(a_i,b_i)$:
a difference of integer slopes dominates the fixed offset difference.
Equal triples have already been collected. Lemma~\ref{m:lem:head-window}
and the remainder estimate in the balanced proof therefore show that
the sign is the coefficient of the greatest triple
$(\lambda_i,a_i,b_i)$. This use of the balanced \emph{numerical estimate}
requires no certificate or comparison algorithm for $E$ itself.
If no heads remain, use $T=\max(1,|K|+1)$; the final sign of $rE+K$ is
$\sgn(r)$ when $r\ne0$, otherwise $\sgn(K)$. Denote the computed final
sign by $\sigma$.

Let $V$ be the actual raw support of $E$. Empty support is handled by one
closed evaluation. Otherwise, outside $[1,T-1]^V$ there is an $x\ge T$,
and positivity gives $E\ge T$. Evaluate every row in the finite box exactly
and record its sign. For a desired sign $s$ the exact formula is
\[
 \left(\bigvee_{x\in V}x\ge T\right)\wedge[\sigma=s]
 \quad\vee\quad
 \bigvee_{\substack{v\in[1,T-1]^V\\\sgn Q(v)=s}}
       \bigwedge_{x\in V}x=v(x).
\]
It is a finite Presburger formula on the original variables. The box is
effective and each closed evaluation terminates. If the tail falsifies a
universal comparison, a concrete admissible countervaluation sets one
supported variable to $T$ and all others to one; the dominance property
certifies it without evaluating a giant collapse value.
\end{proof}

For example, even a core $E=h_\alpha(x)+h_\beta(h_\gamma(y)+z)$ outside
a proposed balanced certificate may be used in
$3h_\lambda(E+c)-2h_\mu(2E+d)-E$. Its support is still $(x,y,z)$.
This result is a comparison theorem for that row, not a graph definition
or a general entry-map algorithm for the core. An unrelated open term
outside $E$, a variable ordinal head, or a nonlinear dependence on $E$
is outside the statement.


\section{Consequences for literal canonical assignments}
\label{sec:source-applications}

The numerical results apply to original, undecorated G\"odel-$T$ terms.
Fix one computable admissible policy for expanding the finite sums and
repetitions in the assignment clauses.  The canonical assignment here is
the resulting zero-history assignment of the literal input term.
Normalizing the term first would define a different object.  The type
level is
\[
 \lv(o)=0,\qquad
 \lv(\sigma\to\tau)=\max\{\lv(\sigma)+1,\lv(\tau)\}.
\]
Vectors are padded by zeros when their levels differ.  At any admissible
valuation, weak comparison means weakness in every coordinate;
Wilken--Weiermann strict comparison means strictness in coordinate zero
and weakness in every positive coordinate.  Equality means equality in
every coordinate.  The strict relation is not a lexicographic order.

We use the following exact source interface.  Its domain and construction
are established from the original assignment clauses in
\cite[the canonical-source section]{wwcomparison}; the source notation
and original bounded-norm conditions are those of \cite{ww}.

\begin{proposition}[Canonical forward interface]
\label{prop:source-interface}
Let $A$ be an original well-typed term all of whose free variables have
ground type.  Its literal canonical vector is computable.  Every positive
coordinate is a closed computable ordinal, and its scalar has an
effectively constructed, valuation-preserving expression in
$\mathcal P_d$, where the $d$ inputs are the distinct ground free names.
Repeated names remain one input.  Arbitrary finite internal and bound
types, result types, recursor instances, and nesting are allowed.
Every positive integer is an admissible value of a ground component,
independently for distinct ground names.
\end{proposition}

The relevant invariant explains the interface's scope.  An actual positive
coordinate is generated from $1,\omega$, positive-index component
variables, natural sums, and graded terms
$2^a\otimes b$, where the exponent comes from the next coordinate.
It contains no promotion of a scalar collapse to a positive coordinate.
The box operation, partials and abstraction preserve this invariant
simultaneously.  Abstraction removes its bound component variables;
ground free names have no positive components.  Hence positive outputs
are closed.  An actual scalar contains only its scalar variables,
positive constants, addition, fixed finite repetitions and
\[
 \operatorname{Psi}(c,u)=\psi(\omega\otimes c+u).
\]
Its positive head $c$ is closed, and $\omega\otimes c$ is a fixed
nonzero limit.  The syntactic abstraction multiplier is fixed before
valuation.  These facts give precisely the stated positive expression
and preserve correlations.  The longer domain proof, including the
correct weak scalar partial estimate and the original abstraction
budget, is the cited source theorem; membership in an unrestricted
ambient ordinal-expression grammar would not suffice.

\subsection{Equality and one shared ground name}

\begin{corollary}[Universal source equality]
\label{cor:source-equality}
Universal equality of literal canonical vectors is decidable for every
original source pair with ground free variables, with no bound on its
internal types.
\end{corollary}
\begin{proof}
Compute the two vectors and compare their finitely many closed positive
coordinates.  If they agree, apply the complete normal form of
Theorem~\ref{u:thm:identity} to the two scalar expressions.
Different forms give the computed positive countervaluation from that
theorem.  The comparison substitutes admissible ground coordinates;
it does not substitute source numerals and recompute a different
canonical assignment.
\end{proof}

\begin{corollary}[One-name Boolean comparison]
\label{cor:source-one-name}
Take a finite Boolean combination of pointwise canonical equality, weak
comparison, and Wilken--Weiermann strict comparison atoms.  If the union
of its free names is contained in one ground name $X:o$, either its
existential or universal closure over the admissible component of $X$
is decidable.  Internal types and result levels are unrestricted.
\end{corollary}
\begin{proof}
For an atom comparing $A$ and $B$, compute its positive-coordinate guard
\[
 U_{\le}=\bigwedge_{i>0}\can A_i\le\can B_i,
 \qquad
 U_=\bigwedge_{i>0}\can A_i=\can B_i,
\]
using zero padding.  The three pointwise atoms are respectively
\[
 U_=\wedge(E_A=E_B),\qquad
 U_{\le}\wedge(E_A\le E_B),\qquad
 U_{\le}\wedge(E_A<E_B).
\]
Place each computed guard inside its original Boolean atom and apply
Corollary~\ref{u:cor:boolean}.  In particular a false guard is not
removed from beneath a negation.  The admissible ground domain is exactly
$\NN$ by Proposition~\ref{prop:source-interface}.
\end{proof}

One name in each term is not sufficient if the two names differ: the
one-input hypothesis concerns the whole query's shared input set.

\subsection{Finite coefficients and recursor-free terms}

\begin{theorem}[The finite-coefficient class]
\label{thm:source-finite}
For a ground-open source pair, suppose every literal scalar collapse
$\psi(\omega\otimes c+g)$ in the generated assignments has a finite
closed coefficient $c$.  Universal weak comparison, Wilken--Weiermann
strict comparison, and equality are decidable for the exact standard
control.  Every recursor-free ground-open source satisfies this condition,
regardless of its internal type levels.
\end{theorem}
\begin{proof}
The hypothesis is decidable by finite traversal of the actual scalar
trees and closed ordinal comparison with $\omega$.  For $c=k<\omega$,
Corollary~\ref{f:cor:finite-coefficient} replaces the call exactly by
\[
          g\longmapsto f^k(g+2k)+k.
\]
The iterate length is a closed integer determined from the source.  Its
expansion is therefore a finite expression in $(\N;0,1,+,<,f)$,
not a variable-iteration primitive.  Retain every original name and apply
Theorem~\ref{m:thm:snyatkov} to the required universal scalar formula,
with $x_i\ge1$.  The closed positive-coordinate tests complete the
corresponding vector comparison.

For the recursor-free assertion, inspect the strengthened positive
coordinate invariant in Proposition~\ref{prop:source-interface} before
erasing the bound components.  The only source constant introducing
$\omega$ into a positive coordinate is a recursor.  Without recursors,
the positive construction uses $1$, component variables, natural sums
and products, and $2^a$.  After the ground-free restriction, all bound
components have been replaced or removed by the literal partials and
all-ones substitutions.  Induction through those same operations leaves
finite positive values: sums, products, and $2^n$ preserve finiteness.
This holds at every coordinate and every internal type.  Thus every
direct scalar coefficient is finite.
\end{proof}

A recursor is not itself a certificate of failure.  For example the
canonical vector of $R^oX$ is $(f(x+4)+1,1,\omega)$; its displayed
scalar collapse has coefficient one.  The theorem checks the actual
coefficient condition rather than a syntactic ban on all recursors.

\subsection{Polynomial heads and internal type level two}

\begin{lemma}[Polynomial direct heads]
\label{lem:source-polynomial}
If every internal type and subtype in a ground-open source has level at
most two, all its direct scalar heads are polynomial limits below
$\omega^\omega$.
\end{lemma}
\begin{proof}
Follow the same generated positive-coordinate invariant.  There is no
coordinate three available.  At coordinate two an exponent constructor
would require such a coordinate, so its closed values are natural linear
combinations of $1$ and $\omega$.  The positive box and abstraction
clauses only add such values and take fixed natural multiples there.
At coordinate one the exponential factors therefore have the form
\[
                  2^{a\omega+b}=\omega^a2^b,
                  \qquad a,b<\omega.
\]
Natural sums and products of these factors are polynomials in $\omega$.
Whole-free partial copies and fixed repetitions preserve this property;
temporary component variables disappear in abstraction.  It follows
that every closed coefficient $c$ of a direct scalar collapse is a
polynomial in $\omega$, so $\omega\otimes c$ is a polynomial limit
below $\omega^\omega$.

The hypothesis includes the full types of recursor constants and of
temporary binders, not just the result type or the free context.  A
higher-level recursor could introduce a coordinate above two and would
invalidate this induction even if the final result were ground.
\end{proof}

\begin{theorem}[Comparison through internal level two]
\label{thm:source-leveltwo}
Universal weak comparison, Wilken--Weiermann strict comparison, and
equality of literal canonical assignments are decidable for ground-open
sources whose maximum internal type level is at most two.
\end{theorem}
\begin{proof}
Compute the closed positive-coordinate tests and the scalar circuit from
Proposition~\ref{prop:source-interface}.  Its heads are polynomial by
Lemma~\ref{lem:source-polynomial}.  Translate all circuit nodes by the
functional graphs in Theorem~\ref{m:thm:graphs}.  Quantify the fresh
natural graph variables together with the original positive inputs and
conjoin the graph equations positively.  The counterexample output test
is $E_A>E_B$ for weak comparison, $E_A\ge E_B$ for strict comparison,
and $E_A\ne E_B$ for equality.  Thus no graph equation is negated when
the requested relation is negated.

Lemma~\ref{m:lem:graph-domains} splits the finite graph exceptions and
puts each remaining case on its individual word domains.
Theorem~\ref{m:thm:existential} decides the resulting closed existential
sentences.  A counterexample exists exactly when the corresponding
universal scalar relation fails.  The needed degree bound is finite and
computed from the input pair; there is no fixed degree bound across all
terms of internal level two.
\end{proof}

At arbitrary internal levels the common-head, balanced, or single-core
theorems can still apply to a scalar residual carrying the corresponding
finite certificate.  One first computes the literal assignment and then
checks the certificate.  Failure of a proposed certificate is not a
negative answer to comparison.  No reverse source-realization theorem
for every auxiliary numerical expression is needed in any of these
applications.


\section{The remaining standard-control boundary}

The normal form decides universal identity at arbitrary fixed heads and
in arbitrary finite arity.  It does not decide whether a nonidentity has
some multivariate zero.  The one-input algorithm does decide every
one-block Boolean comparison, because its truth is constant beyond a
computed bound.  Its proof uses a common growing input: a peak depending
only on a small coordinate need not dominate an affine term in a different
large coordinate.

Polynomial heads permit independent raw inputs through their exact
finite-degree kernels and counts.  The elimination proved here is a
closed-sentence induction on the number of quantified raw variables.
It does not preserve arbitrary free parameters through that induction,
and therefore does not establish the full alternating first-order theory
of polynomial-collapse functions.  The common-head and balanced theorems
have a different scope: their actual parameter-preserving formulas allow
arbitrary quantifier nesting by reduction to the standard-control base
theory.

Unrestricted order comparison of multivariate positive expressions with
arbitrary nonpolynomial heads and no such certificate is not decided by
these methods.  This is a question about the unchanged standard control,
separate from the all-ground decision theorem for the permitted controls
constructed in \cite{wwcomparison}.  The exact kernels and sparse
normal forms remain useful for that standard-control question and for
the structure of the resulting numerical functions.  Source realization
and the effects of varying a control are studied further in
\cite{wwsource}.

\section*{Preparation and verification status}

This manuscript consolidates earlier work on positive identities,
one-input comparison, and multivariate fragments.  It was prepared through
AI-assisted research directed by the repository owner; this preparation
record assigns no individual human mathematical authorship.  The paper
contains hand proofs.  Earlier formalization checked many exact source
interfaces and exposed a predecessor-convention error that is corrected
in the kernel definition here.  The complete closed existential assembly,
its source consumers, and the published first-order arithmetic supplier
are not claimed to have been verified together as one proof-assistant
development.  Build checks and separate AI review do not constitute
independent human verification or a historical priority assessment.

\begin{thebibliography}{99}

\bibitem{ww}
Gunnar Wilken and Andreas Weiermann.
Derivation lengths classification of G\"odel's T extending Howard's assignment.
\emph{Logical Methods in Computer Science} \textbf{8}(1:19) (2012), 1--44.
\href{https://doi.org/10.2168/LMCS-8(1:19)2012}{doi:10.2168/LMCS-8(1:19)2012}.
\url{https://lmcs.episciences.org/1073}.

\bibitem{snyatkov}
A. S. Snyatkov.
Decidability of the theory
$T_f=\operatorname{Th}(\omega,0,1,<,+,f(x),F(x))$.
\emph{Vestnik Tverskogo Gosudarstvennogo Universiteta,
Series Applied Mathematics}, no. 2(9) (2008), 39--51.  In Russian.
\url{https://eprints.tversu.ru/id/eprint/494/1/19950136_2008_2_snyatkov.pdf}.

\bibitem{cooper}
D. C. Cooper. Theorem proving in arithmetic without multiplication.
In B. Meltzer and D. Michie, editors, \emph{Machine Intelligence 7},
91--99. Edinburgh University Press, 1972.
\url{https://www21.in.tum.de/teaching/logik/SS16/Exercises/Cooper.pdf}.

\bibitem{wwcomparison}
\emph{Decidable Ground Comparison for Norm-Controlled Ordinal Assignments}.
Research manuscript, 2026.
\href{https://github.com/avg-netizen/logic-lineages-papers/blob/papers/ii5-explanatory-2026-11-17/17-permitted-control-canonical-comparison/paper.pdf}{Companion manuscript}.

\bibitem{wwsource}
\emph{Source Expressivity and Control Sensitivity of Ordinal Assignments}.
Research manuscript, 2026.
\href{https://github.com/avg-netizen/logic-lineages-papers/blob/papers/ii5-explanatory-2026-11-17/19-source-expressivity-and-control-sensitivity/paper.pdf}{Companion manuscript}.

\bibitem{wwcounts}
\emph{Controlled Ordinal Descents, Restricted Partitions and Finite Cartier Observers}.
Research manuscript, 2026.
\href{https://github.com/avg-netizen/logic-lineages-papers/blob/papers/ii5-explanatory-2026-11-17/20-controlled-descents-and-cartier-observers/paper.pdf}{Companion manuscript}.

\end{thebibliography}

\end{document}
